\providecommand{\hecsgn}{\operatorname{sgn}}
\providecommand{\hecVar}{\operatorname{Var}}
\providecommand{\hecCov}{\operatorname{Cov}}
\providecommand{\hecE}{\mathbb{E}}
\providecommand{\hecind}{\mathbf{1}}
\providecommand{\nabFix}{\operatorname{Fix}}
\providecommand{\nabIrr}{\operatorname{Irr}}
\providecommand{\nabsgn}{\operatorname{sgn}}
\providecommand{\nabone}{\mathbf{1}}
\providecommand{\nabpos}[1]{(#1)^{+}}
\section{Detectors for specific Galois groups and fields}\label{sec:grp}

\providecommand{\grpIrr}{\operatorname{Irr}}
\providecommand{\grpSt}{\mathrm{St}}
\providecommand{\grpFix}{\operatorname{Fix}}
\providecommand{\grpsgn}{\operatorname{sgn}}
\providecommand{\grpstd}{\mathrm{std}}
\providecommand{\grpone}{\mathbf{1}}
\providecommand{\grpDiff}{\mathfrak{D}}
\providecommand{\grprd}{\operatorname{rd}}
\providecommand{\grppos}[1]{#1^{+}}
% Long file names and identifiers set in typewriter cannot be hyphenated; the
% group opened here (closed at the end of the section) lets TeX break before them.
\begingroup\emergencystretch=3em

The general reductions of \cref{sec:nab} turn the Riemann hypothesis for the
Dedekind zeta function of a Galois extension into a one-sided inequality for a
sum over rational primes, weighted by a class function of the Frobenius
element. This section records what is gained for particular groups and fields:
explicit character identities which certify holomorphy and multiplicity
domination without Artin's holomorphy conjecture, statements of optimality
within declared classes of detectors, and the resulting inequality for primes.
The finite group theory and the comparisons of zero orders are unconditional.
The prime inequalities are not proved: every criterion below is a reduction,
and no Riemann hypothesis beyond \cref{src:thm:dgrh}, and no case of Artin's
primitive-root conjecture, is established.

\subsection{Conventions}\label{grp:sec:conv}
{\rightskip=0pt plus 1fil\relax
\sources{grh_a6_transfer_2026-11-25; grh_sylow_transfer_2026-11-25; psl2_grh_transfer; grh_deficit_transfer_2026-11-25; grh-transfer}
}

Unless the contrary is stated, $L/\Q$ is a finite Galois extension with group
$G$. The analytic interface is \cref{nab:sec:interface}
(\cref{grp:prop:interface}, \cref{grp:rem:variants}, \cref{grp:cor:detector}), stated
there for a Galois field $E$, which is $L$ here: the ramp $P_w$, its centred remainder
$R_w$, the sums $\theta_w$, the mean $\delta_w$ and the prime Dirichlet series
$D^{\mathrm{pr}}_w$ are those defined there, with $\Frob_p$ for $g_p$. For a subgroup $H\le G$ we write $L^H$ for its fixed field and
$\pi_H=\Ind_H^G\triv$ for the permutation character on $G/H$, so that
\[
 \pi_H(g)=\#\grpFix(g,G/H)=\frac{\abs{\set{x\in G:x^{-1}gx\in H}}}{\abs{H}},
 \qquad
 (\Ind_H^G\lambda)(g)=\frac{\abs{C_G(g)}}{\abs{H}}\sum_{h\in H\cap g^G}\lambda(h)
\]
for a linear character $\lambda$ of $H$. We use Artin formalism in its
classical form \cite{Neukirch1999Algebraic,Booker2006Artin}: the Artin
$L$-function is additive in the character and compatible with induction,
$L(s,\Reg_G)=\zeta_L(s)$ and $L(s,\pi_H)=\zeta_{L^H}(s)$, all local factors
included. If $\lambda$ is a nontrivial linear character of $H$, then
$L(s,\Ind_H^G\lambda)$ is the Hecke $L$-function of a nontrivial finite-order
character of $L^H$; it is entire and does not vanish on $\Re s\ge1$. By
Brauer's theorem every Artin $L$-function is meromorphic. For $\rho\in\C$,
$\ord_\rho$ denotes the signed order, positive at a zero and negative at a
pole. No holomorphy of an individual nonmonomial Artin $L$-function is assumed
anywhere in this section.

For $\rho\ne1$ the Heilbronn character is the virtual character
\[
 \Theta_\rho=\sum_{\psi\in\grpIrr(G)}\ord_\rho L(s,\psi,L/\Q)\,\psi,
 \qquad \Theta_\rho(1)=\ord_\rho\zeta_L .
\]
By \cref{nab:lem:heilbronn}(a) \cite{FooteGinsbergMurty2015Heilbronn}, for every
subgroup $C\le G$ and every $\eta\in\grpIrr(C)$,
\begin{equation}\label{grp:eq:heil}
 \langle\Res_C\Theta_\rho,\eta\rangle=\ord_\rho L(s,\eta,L/L^C),
\end{equation}
and consequently $\Res_C\Theta_\rho$ is a genuine character whenever every
irreducible character of $C$ is monomial, in particular for $C$ cyclic or of
prime-power order. Hence $d=\ord_\rho\zeta_L\ge0$, $\abs{\Theta_\rho(g)}\le d$
for all $g\in G$, and $d=0$ forces $\Theta_\rho=0$: at a point $\rho\ne1$ where
$\zeta_L$ does not vanish, every Artin $L$-function of $L/\Q$ is holomorphic
and nonzero.

\subsection{Icosahedral fields}\label{grp:sec:a5}
{\rightskip=0pt plus 1fil\relax
\sources{grh_method_transfers_2026-11-25_1; grh_sharp_transfers_2026-11-25; psl2_grh_transfer}
}

Let $G=\Gal(L/\Q)\simeq A_5$. Write $\chi,\chi',\eta,\nu$ for the irreducible
characters of degrees $3,3,4,5$, and $\phi=(1+\sqrt5)/2$, $\phi'=(1-\sqrt5)/2$.
Choose subgroups $A_4$, $D_5$, $C_5$, $V_4$ of orders $12,10,5,4$ with
$V_4\le A_4$ and $C_5\le D_5$, and put
\[
 K_5=L^{A_4},\qquad K_6=L^{D_5},\qquad K_{12}=L^{C_5},
\]
of degrees $5,6,12$, with permutation characters $\pi_5,\pi_6,\pi_{12}$.
\begin{center}
\begin{tabular}{c|rrrrr}
\toprule
class&$1$&$2A$&$3A$&$5A$&$5B$\\
size&$1$&$15$&$20$&$12$&$12$\\\midrule
$\chi$&$3$&$-1$&$0$&$\phi$&$\phi'$\\
$\chi'$&$3$&$-1$&$0$&$\phi'$&$\phi$\\
$\eta$&$4$&$0$&$1$&$-1$&$-1$\\
$\nu$&$5$&$1$&$-1$&$0$&$0$\\\midrule
$\pi_5$&$5$&$1$&$2$&$0$&$0$\\
$\pi_6$&$6$&$2$&$0$&$1$&$1$\\
$\pi_{12}$&$12$&$0$&$0$&$2$&$2$\\\midrule
$W_{16}=\pi_5+\pi_{12}-\triv$&$16$&$0$&$1$&$1$&$1$\\
$W_{21}=\pi_5+\pi_6+\pi_{12}-2\,\triv$&$21$&$1$&$0$&$1$&$1$\\
$F=\chi^4$&$81$&$1$&$0$&$\phi^4$&$\phi'^4$\\
\bottomrule
\end{tabular}
\end{center}
Three detectors are available. The characters $W_{16}$ and $W_{21}$ are
nonnegative with mean $1$; $W_{16}$ is the case $q=4$ of the even
$\PSL_2(q)$ detector of \cref{grp:sec:psl-even} and has the least possible
degree, while $W_{21}$ is the detector built from the three subfields
$K_5,K_6,K_{12}$. The tensor character $F=\chi^4$ has degree $81$ and mean
$3$; it does not compete in degree, and its purpose is different: it is the
first tensor power of a three-dimensional representation whose holomorphy and
zero detection are certified by proper-subgroup information.

\subsubsection{Two positive permutation detectors}

\begin{lemma}[Character identities; \stU]\label[lemma]{grp:lem:a5-identities}
Let $\lambda_0$ be a nontrivial character of $V_4$ and $\omega$ a nontrivial
character of $A_4/V_4$. Then
\begin{align}
 \Reg_{A_5}&=4\pi_5+2\pi_6+3\pi_{12}-8\,\triv, \label{grp:eq:a5-reg}\\
 W_{16}&=\triv+\Ind_{V_4}^{A_5}\lambda_0=\triv+\chi+\chi'+\eta+\nu,
 \label{grp:eq:a5-w16}\\
 W_{21}&=W_{16}+\nu,\qquad \nu=\pi_6-\triv=\Ind_{A_4}^{A_5}\omega .
 \label{grp:eq:a5-w21}
\end{align}
Consequently
\begin{gather}
 \zeta_L(s)\zeta(s)^8=\zeta_{K_5}(s)^4\zeta_{K_6}(s)^2\zeta_{K_{12}}(s)^3,
 \label{grp:eq:a5-zeta}\\
 L(s,W_{16})=\frac{\zeta_{K_5}(s)\zeta_{K_{12}}(s)}{\zeta(s)}
 =\zeta(s)L(s,\lambda_0,L/L^{V_4}),\qquad
 L(s,W_{21})=\frac{\zeta_{K_5}(s)\zeta_{K_6}(s)\zeta_{K_{12}}(s)}{\zeta(s)^2}.
 \notag
\end{gather}
\end{lemma}

\begin{proof}
The permutation rows follow from the coset formula, and
\eqref{grp:eq:a5-reg} holds in each of the five columns. The induction formula
gives $\Ind_{V_4}^{A_5}\lambda_0=(15,-1,0,0,0)$, since the three involutions
of $V_4$ have $\lambda_0$-sum $-1$ and centralizer of order $4$, and
$\Ind_{A_4}^{A_5}\omega=(5,1,-1,0,0)$, since $\omega$ sums to $3$ on the
involutions and to $-4$ on the eight elements of order three of $A_4$. The
remaining equalities are read from the table. Artin formalism, applied first
in $\Re s>1$ and then continued meromorphically, gives the identities of
$L$-functions; it does not assert holomorphy of individual nontrivial Artin
factors \cite{GunHazraSahu2025Holomorphy,LemkeOliverThornerZaman2024Approximate}.
\end{proof}

Both detectors are monomially certified in the sense of
\cref{grp:cor:detector}, with $a=1$. Their coefficients at all prime powers
are nonnegative, $0\le a_J(n)\le J(1)\Lambda(n)$, because each summand
$J(g_p^kh)$ in \cref{grp:rem:variants} is nonnegative; at an unramified prime
$a_J(p^k)=J(g_p^k)\log p$, and using only $J(g_p)$ at all powers would be
incorrect. For example the local factors of $L(s,W_{21})$ at unramified
primes, with $t=p^{-s}$, are $(1-t)^{-21}$ for the identity,
$(1-t)^{-1}(1-t^2)^{-10}$ for $2A$, $(1-t^3)^{-7}$ for $3A$ and
$(1-t)^{-1}(1-t^5)^{-4}$ for $5A,5B$. In terms of Euler--Kronecker constants,
the centring constant of \cref{grp:rem:variants} for $W_{21}$ is
$c=-1-\gamma_{K_5}-\gamma_{K_6}-\gamma_{K_{12}}+2\gamma$.

\begin{theorem}[Positive $A_5$ detectors; \stR]\label{grp:thm:a5-criterion}
Use the Riemann hypothesis for $\zeta$ from \cref{src:thm:dgrh}. Suppose that for
every $\eps>0$ there are $C_\eps,X_\eps$ with
\begin{equation}\label{grp:eq:a5-bound}
 R_{W_{21}}(x)\le C_\eps x^{-1/2+\eps}\qquad(x\ge X_\eps).
\end{equation}
Then $\zeta_L$, $\zeta_{K_5}$, $\zeta_{K_6}$ and $\zeta_{K_{12}}$ satisfy the
Riemann hypothesis, and every Artin $L$-function of $L/\Q$ is holomorphic and
nonzero in $1/2<\Re s<1$. Conversely, the Riemann hypothesis for $\zeta_L$
implies $R_{W_{21}}(x)=O(x^{-1/2}\log^2(2x))$. In prime-only form the weights are
$21,1,0,1,1$ on the classes $1,2A,3A,5A,5B$, with main term $\log x$.
\end{theorem}

The detector $W_{16}$, with the weights $16,0,1,1,1$, is by \eqref{grp:eq:a5-w16} the
case $q=4$ of the detector $W$ of \cref{grp:sec:psl-even}, with $B=A_4$, $T'=C_5$ and
$U=V_4$: \cref{grp:thm:psl-transfer,grp:cor:psl-root} give the same equivalence for it,
and \eqref{grp:eq:psl-ineq2} reads
$3\ord_\rho L(s,W_{16})\le\ord_\rho\zeta_L\le5\ord_\rho L(s,W_{16})$ for
$1/2<\Re\rho<1$.

\begin{proof}[Proof of \cref{grp:thm:a5-criterion}]
Fix $1/2<\beta<1$. By \eqref{grp:eq:a5-bound} and
\cref{grp:cor:detector}(i), $L(s,W_{21})$ has no zero in $\Re s>\beta$. The
Riemann hypothesis for $\zeta$ ensures that the denominator $\zeta(s)^2$ has
no zero in $\Re s>1/2$, and the three numerator factors are holomorphic away
from $1$. In that half-plane, away from $1$, the order of $L(s,W_{21})$ is
therefore the sum of the three nonnegative orders of the numerator zeta
functions, whose zeros are excluded individually. Equation
\eqref{grp:eq:a5-zeta} then excludes zeros of $\zeta_L$. Let
$\beta\downarrow1/2$ and use the four functional equations. The statement
about Artin $L$-functions is \eqref{grp:eq:heil}, and the converse is
\cref{grp:prop:interface}(d).
\end{proof}

\begin{remark}[How the two detectors differ]\label[remark]{grp:rem:a5-compare}
By \eqref{grp:eq:a5-w21}, $R_{W_{21}}=R_{W_{16}}+R_\nu$, and
$L(s,W_{21})=L(s,W_{16})\,L(s,\omega,L/L^{A_4})$. The detector $W_{16}$ ignores
the involution class and sees the elements of order three; $W_{21}$ ignores
the class $3A$ and sees the involutions. Neither the inequality
\eqref{grp:eq:a5-bound} nor its analogue for $W_{16}$ formally implies the other,
although each is equivalent to the Riemann hypothesis for $\zeta_L$.

By \cref{grp:thm:psl-detector}, $W_{16}$ is the unique nonnegative genuine
character of least degree containing every nontrivial irreducible character
of $A_5$. The optimality of $W_{21}$ is of a narrower kind. In the normalized
three-field family
\[
 W_{a,b,c}=a\pi_5+b\pi_6+c\pi_{12}+(1-a-b-c)\triv,\qquad a,b,c\ge1,
\]
whose members have mean $1$ and detect each of the three Dedekind factors with
coefficient at least $1$, the identity value is
$W_{a,b,c}(1)=1+4a+5b+11c\ge21$; so every nonnegative member has peak at least
$21$, attained by $a=b=c=1$, also for real coefficients. The constraint
$b\ge1$ excludes $W_{16}=W_{1,0,1}$. The direct regular-character detector has
peak $60$ at the same mean.
\end{remark}

\paragraph{Sparse weights of both signs.}
For a weight $0\le f\le1$ with $f(1)=0$ whose character coefficients may have both
signs, \cref{nab:prop:kappa} gives $P_+(f)\ge\kappa(G)P_-(f)$ as soon as $f$ vanishes on
one nonidentity class.
For $A_5$ the character table gives the ratios $3/5$, $5/8$ and $(10-2\sqrt5)/15$ at
elements of orders $2$, $3$ and $5$: at one class of $5$-cycles the extreme normalized
traces are $(1+\sqrt5)/6$ for a representation of degree three and $-1/4$ for the
representation of degree four. So omitting even one nonidentity class costs at least
$(10-2\sqrt5)/15$ in the ratio $P_+/P_-$.

If $f$ on $A_5$ is required to depend only on the order of the Frobenius element, that is
on the cycle type in the natural action of degree five, write $a,b,c\ge0$ for its values at
orders $2,3,5$. With class sizes $(1,15,20,12,12)$, and the two characters of degree three
combined, the signed contributions $60\chi(1)a_\chi$ to the masses are
\[
 v_3=-90a+72c,\qquad v_4=80b-96c,\qquad v_5=75a-100b ,
\]
and $60P_\pm$ is the sum of the positive, respectively negative, parts of these numbers.
Since $\sum_iv_i=-15a-20b-24c$, every nonzero admissible weight has $P_->0$.

\begin{proposition}[The rational frontier for $A_5$; \stU]\label{nab:prop:a5frontier}
The minimum of $P_+/P_-$ over nonzero weights $(a,b,c)$ supported on at most $B$ group
elements is:
\begin{center}
\begin{tabular}{c|c|c}
$B$ & $\min P_+/P_-$ & $(a,b,c)$ attaining it\\\hline
$0$--$14$ & no nonzero weight & --\\
$15$--$19$ & $5/6$ & $(1,0,0)$\\
$20$--$23$ & $4/5$ & $(0,1,0)$\\
$24$--$34$ & $3/4$ & $(0,0,1)$\\
$35$--$38$ & $2/3$ & $(1,3/4,0)$\\
$39$--$43$ & $5/8$ & $(4/5,0,1)$\\
$44$--$58$ & $3/5$ & $(0,1,5/6)$\\
$59$--$60$ & $0$ & $(1,1,1)$
\end{tabular}
\end{center}
\end{proposition}
\begin{proof}
Supporting one order fixes the ratio. Below budget $59$ at most two orders are supported;
normalize their values to sum to one. On each interval where the signs of the $v_i$ are
fixed the ratio of linear forms is monotone or constant, so its minimum lies at an axis or
at a zero of some $v_i$. The three interior candidate rays are $(4,3,0)$, $(4,0,5)$ and
$(0,6,5)$, which give the table by substitution; at budget $59$ the full indicator has
$P_+=0$. The frontier is recomputed in exact arithmetic by
\artifact{grh_sharp_torsion_and_symplectic_blocks/checks/check_support.py}.
\end{proof}

In particular, monitoring fewer than all three nonidentity splitting types of an $A_5$
quintic forces a positive contribution of at least three fifths of the negative mass. This
is a statement about the spectrum of the observable; converting mixed signs into a
conclusion on zeros would need additional information about their multiplicities.

\begin{corollary}[Mean-square distance from congruence weights; \stU]
\label[corollary]{grp:cor:a5-variance}
For every fixed modulus $q$ and every $a:(\Z/q\Z)^\times\to\C$, in the notation
of \cref{nab:thm:periodic},
\[
 \mathcal E_q(W_{21},a)=7+\frac1{\varphi(q)}\sum_{b\in(\Z/q\Z)^\times}\abs{a(b)-1}^2,
 \qquad
 \mathcal E_q(W_{16},a)=4+\frac1{\varphi(q)}\sum_{b\in(\Z/q\Z)^\times}\abs{a(b)-1}^2 .
\]
\end{corollary}

This is \cref{nab:thm:periodic} for the perfect group $A_5$, where $H_q=G$ and
$m_q=1$: by \eqref{nab:eq:energy} the minimum is $1+1+1+2^2=7$ for
$W_{21}=\triv+\chi+\chi'+\eta+2\nu$ and $4$ for $W_{16}$. It does not cover moduli
changing with the scale of the prime sum (\cref{nab:rem:fixed}).

\begin{example}[A specified field]\label[example]{grp:ex:a5-field}
One can take $L$ to be the splitting field of $P(X)=X^5+20X+16$. The
discriminant is $1\,024\,000\,000=32\,000^2$, so the Galois group lies in
$A_5$. Modulo $3$, $P=X^5-X+1$ has no root, and its remainders modulo the
three monic irreducible quadratics $X^2+1$, $X^2+X+2$, $X^2+2X+2$ are $1$,
$X+1$, $X+1$; hence $P$ is irreducible modulo $3$. Modulo $7$,
$P=(X+2)(X+3)(X^3+2X^2-2X-2)$ with the cubic having no root. Thus the group
contains elements of orders $5$ and $3$; its order is $15$, $30$ or $60$. A
group of order $15$ is cyclic and does not embed in $S_5$, and $A_5$, being
generated by elements of order three, has no subgroup of index two. Hence the
group is $A_5$. The certificate is reproduced by
\artifact{grh_method_transfers_2026-11-25_1/code/a5_detector.py}.
\end{example}

\subsubsection{The faithful tensor construction}

\begin{proposition}[Tensor amplification; \stU]\label[proposition]{grp:prop:tensor}
Let $L/K$ be a finite Galois extension of number fields with nontrivial group
$G$, let $\varrho$ be a faithful complex representation of dimension $r$, and
put $\tau=\triv\oplus\varrho$, $D=r+1$,
\[
 F_m(g)=\abs{\chi_\tau(g)}^{2m},\qquad
 q_\tau=\max_{g\ne1}\frac{\abs{\chi_\tau(g)}^2}{D^2},\qquad
 \kappa_m=\frac{F_m(1)-\sum_{g\ne1}F_m(g)}{\abs G}.
\]
Then $q_\tau<1$, and $F_m$ is the pointwise nonnegative character of
$(\tau\otimes\overline\tau)^{\otimes m}$. If $(\abs G-1)q_\tau^m<1$, then
$\kappa_m>0$ and for every $s_0\ne1$, with $d=\ord_{s_0}\zeta_L$,
\begin{equation}\label{grp:eq:tensor-domination}
 \kappa_m\,d\le\ord_{s_0}L(s,F_m,L/K)\le\langle F_m,\triv\rangle\,d .
\end{equation}
Consequently $L(s,F_m)$ is holomorphic away from $1$, has there the same zero
support as $\zeta_L$, and has a pole of order $\langle F_m,\triv\rangle$ at
$1$. If $e$ is the exponent of $G$, then
$q_\tau\le1-4\sin^2(\pi/e)/D^2$.
\end{proposition}

\begin{proof}
If $\abs{\chi_\tau(g)}=D$, equality in the triangle inequality forces every
eigenvalue of $\tau(g)$ to equal the distinguished eigenvalue $1$, so $g=1$ by
faithfulness. Let $\Theta=\Theta_{s_0}$ be the Heilbronn character of $L/K$;
then $\abs{\Theta(g)}\le d$ by \eqref{grp:eq:heil}. The inner product
$\langle F_m,\Theta\rangle=\ord_{s_0}L(s,F_m)$ is real; its identity term is
$F_m(1)d/\abs G$ and every other term has absolute value at most
$F_m(g)d/\abs G$. This gives both bounds. The pole order at $1$ is the
multiplicity of the trivial character, the nontrivial Artin $L$-functions
being holomorphic and nonzero at $1$. For the last bound diagonalize
$\varrho(g)$ unitarily: $D^2-\abs{\chi_\tau(g)}^2$ is the sum of the squared
pairwise distances among the $D$ eigenvalues, and for $g\ne1$ the pair
consisting of $1$ and an eigenvalue $\lambda\ne1$, an $e$th root of unity,
contributes at least $4\sin^2(\pi/e)$.
\end{proof}

The proposition establishes holomorphy for an aggregate tensor representation,
not for the constituents of $\varrho$. Tensor generation itself is classical;
compare \cite{Steinberg2014Burnside}. The exponent bound is crude, and for
$A_5$ a character-table certificate does better.

\begin{theorem}[Fourth-tensor certificate for $A_5$; \stU]\label{grp:thm:a5-tensor}
Let $\epsilon$ be the sign character of $D_5$, trivial on rotations, $\lambda$
the nontrivial character of $C_5\le D_5$ corresponding to the eigenvalues of
$\chi$, and $\omega$ a nontrivial character of $A_4/V_4$. Then
\[
 A=\Ind_{D_5}^{A_5}\epsilon=\chi+\chi',\qquad
 B=\Ind_{C_5}^{A_5}\lambda=\chi+\eta+\nu,\qquad
 C=\Ind_{A_4}^{A_5}\omega=\nu,
\]
\[
 \chi^2=\triv+\chi+\nu,\qquad
 \chi^3=\triv+3\chi+\chi'+\eta+2\nu,\qquad
 F=\chi^4=3\,\triv+6\chi+3\chi'+4\eta+7\nu,
\]
and
\begin{equation}\label{grp:eq:a5-fourth}
 4F=\Reg_{A_5}+11\,\triv+9A+12B+11C .
\end{equation}
For every $A_5$-extension $L/\Q$,
\begin{equation}\label{grp:eq:a5-product}
 L(s,F)^4=\zeta_L(s)\zeta(s)^{11}
 L(s,\epsilon,L/L^{D_5})^9L(s,\lambda,L/L^{C_5})^{12}L(s,\omega,L/L^{A_4})^{11},
\end{equation}
$L(s,F)$ is holomorphic on $\C\setminus\set1$ with a pole of order $3$ at $1$,
and at every $s_0\ne1$, with $e_F=\ord_{s_0}L(s,F)$ and $d=\ord_{s_0}\zeta_L$,
\begin{equation}\label{grp:eq:a5-mult}
 d/4\le e_F\le3d .
\end{equation}
\end{theorem}

\begin{proof}
The inductions follow from Frobenius reciprocity with the character table, and
the tensor decompositions by taking inner products; the identity
\eqref{grp:eq:a5-fourth} is then a comparison of the five coefficients, using
$\Reg=\triv+3\chi+3\chi'+4\eta+5\nu$. The exact computation is recorded in
\artifact{grh_sharp_transfers_2026-11-25/data/a5_tensor_certificate.json}.
Artin formalism gives \eqref{grp:eq:a5-product}. The three nontrivial
one-dimensional factors are entire Hecke functions, nonzero at $1$, so the
right side is holomorphic away from $1$ with a pole of order $12$ there. Since
$L(s,F)$ is meromorphic and its fourth power is holomorphic away from $1$, it
is holomorphic there; its pole order at $1$ is the multiplicity $3$ of the
trivial character. Equation \eqref{grp:eq:a5-product} gives $e_F\ge d/4$.
Since $F(g)=\chi(g)^4\ge0$ and $\abs{\Theta_{s_0}(g)}\le d$,
\[
 e_F=\frac1{60}\sum_gF(g)\overline{\Theta_{s_0}(g)}\le\frac d{60}\sum_gF(g)=3d .
 \qedhere
\]
\end{proof}

Thus $L(s,F)$ and $\zeta_L$ have identical zero support, even if individual
Artin factors had poles. This is stronger than observing that every
irreducible character occurs in $F$: positive constituent multiplicities do
not by themselves prevent cancellation of meromorphic orders.

\begin{proposition}[Scope of the optimality; \stU]\label[proposition]{grp:prop:a5-cone}
\begin{enumerate}[label=\textup{(\alph*)}]
\item The virtual character $\Theta_*=-\chi+\chi'+\eta$ has degree $4$, and
its restrictions to the maximal subgroups $A_4$, $D_5$, $S_3$ are genuine:
\[
\begin{array}{c|l}
 A_4&\triv+\text{the three-dimensional irreducible}\\
 D_5&2\cdot\text{one two-dimensional irreducible}\\
 S_3&\triv+\grpsgn+\text{the two-dimensional irreducible}
\end{array}
\]
Moreover $\langle\chi^j,\Theta_*\rangle=-1$ for $j=1,2,3$ and
$\langle\chi^4,\Theta_*\rangle=1$.
\item Consequently the first three tensor powers of $\chi$ cannot be certified
by the proper-subgroup constraints, and no inequality $d\le k\,e_F$ valid
throughout the cone of virtual characters with genuine restrictions to all
proper subgroups can have $k<4$; the identity \eqref{grp:eq:a5-fourth} attains
$k=4$.
\item Write $\Theta=x\chi+y\chi'+z\eta+w\nu$. The restrictions of $\Theta$ to
$A_4,D_5,S_3$ are genuine if and only if
\[
 z\ge0,\quad w\ge0,\quad x+y\ge0,\quad x+z+w\ge0,\quad y+z+w\ge0,
\]
and this cone has the six extreme rays
\[
 (1,0,0,0),\ (0,1,0,0),\ (-1,1,1,0),\ (1,-1,1,0),\ (-1,1,0,1),\ (1,-1,0,1)
\]
in the coordinates $(x,y,z,w)$, on which $e_F/d$ equals
$2,\ 1,\ 1/4,\ 7/4,\ 4/5,\ 2$. Hence the exact bounds in this cone are
$d/4\le e_F\le2d$.
\end{enumerate}
\end{proposition}

\begin{proof}
(a) and (b). The restrictions and inner products are computed from the
character table. Each of the three maximal subgroups is monomial, so by
\eqref{grp:eq:heil} these nonnegative restriction orders are indeed imposed at
an arithmetic point, and every proper subgroup lies in one of them: the
enumeration in
\artifact{grh_sharp_transfers_2026-11-25/code/verify_a5_tensor.py} verifies
the $59$ subgroups of $A_5$ and the maximal subgroup counts, ten of order $6$,
six of order $10$ and five of order $12$. Thus $\Theta_*$ satisfies all
necessary proper-subgroup conditions, including every cyclic restriction,
while $e_F/d=1/4$ on $\Theta_*$ and $\langle\chi^j,\Theta_*\rangle<0$ for
$j\le3$.

(c) The inequalities express the genuineness of the three restrictions. For a
nonzero vector of the cone put $s=x+y\ge0$ and
$t=(x+z+w)/(s+2z+2w)\in[0,1]$. Then the vector is the convex combination, with
coefficients $t$ and $1-t$, of
$s(1,0,0,0)+z(1,-1,1,0)+w(1,-1,0,1)$ and
$s(0,1,0,0)+z(-1,1,1,0)+w(-1,1,0,1)$, which proves that the six rays generate.
The ratios follow from $e_F=6x+3y+4z+7w$ and $d=3x+3y+4z+5w$.
\end{proof}

The optimality is a statement about the algebraic cone of necessary character
conditions. It does not assert that an actual zero of a number field realizes
$\Theta_*$, nor that the first three tensor $L$-functions have poles. Under the
Riemann hypothesis for $\zeta$ the cone bounds of (c), which assume that the
trivial character does not occur, apply throughout $1/2<\Re s_0<1$. The crude
pointwise test of \cref{grp:prop:tensor} fails for $\chi^4$: the identity
value is $81$, while the nonidentity sum is $15+12(\phi^4+\phi'^4)=99$. Among
even powers it first succeeds at $\chi^6$, where
$729>15+12(\phi^6+\phi'^6)=231$. The proper-subgroup certificate therefore
removes one tensor-square amplification step.

\begin{theorem}[Tensor criterion; \stR]\label{grp:thm:a5-target}
Let $L/\Q$ be Galois with group $A_5$ and $F=\chi^4$. The coefficients of
$-L'/L(s,F)$ satisfy
$0\le a_F(p^k)=\log p\,\abs{I_p}^{-1}\sum_{h\in I_p}\chi(g_p^kh)^4\le81\log p$,
including ramified primes. Let $R_F$ be the full prime-power remainder of
\cref{grp:rem:variants}, with $\delta_F=3$. If for every $\eps>0$
\[
 R_F(x)\le C_{L,\eps}x^{-1/2+\eps}\qquad\text{for all sufficiently large }x,
\]
then every nontrivial zero of $\zeta_L$ has real part $1/2$, and every Artin
$L$-function of $L/\Q$ has neither a zero nor a pole off that line inside the
critical strip.
\end{theorem}

\begin{proof}
The coefficient formula is the trace of $g_p^k$ on the inertia invariants of
the fourth tensor power of a representation with character $\chi$. By \cref{grp:thm:a5-tensor}, $L(s,F)$ is holomorphic
away from $1$, and the Mellin transform of $R_F$ has in $\Re s>0$ only simple
poles at the zeros $\alpha$ of $L(s,F)$, with residue
$\ord_\alpha L(s,F)/(\alpha(1-\alpha))$, positive for real $0<\alpha<1$. There
are no zeros with $\Re s\ge1$, by \eqref{grp:eq:a5-product} and the classical
zero-free boundary of its factors. For each $1/2<\theta<1$ the hypothesis and
\cref{ana:thm:onesided} exclude zeros of $L(s,F)$ with real part greater than
$\theta$; the comparison $d/4\le e_F$ excludes zeros of $\zeta_L$ there. Let
$\theta\downarrow1/2$ and use the functional equation. At an off-line point of
the strip $d=0$, and \eqref{grp:eq:heil} forces the whole Heilbronn character
to vanish.
\end{proof}

The last assertion permits poles of individual Artin factors on the critical
line; it does not prove Artin's holomorphy conjecture.

\subsection{The alternating group \texorpdfstring{$A_6$}{A6}}\label{grp:sec:a6}
{\rightskip=0pt plus 1fil\relax
\sources{grh_a6_transfer_2026-11-25}
}

Let $G=A_6$ in its natural action on $\set{1,\dots,6}$; permutations multiply
right to left. Merging the two classes of $5$-cycles, the cycle patterns are
\[
\begin{array}{c|rrrrrr}
 &1^6&1^3\,3&1^2\,2^2&1\,5&2\,4&3^2\\\hline
 \text{size}&1&40&45&144&90&40
\end{array}
\]
The $5$-cycles split into two classes of size $72$, and no other type splits:
a $5$-cycle has only even elements in its centralizer in $S_6$, while every
other type has an odd centralizer element. Hence $G$ has seven conjugacy
classes. All characters constructed below are constant on the six patterns.

\subsubsection{Four subgroups and nine characters}
The following generators specify subgroups and linear characters; an
unmentioned generator has value $+1$.
\begin{itemize}
\item $H_4=\langle v_1,v_2\rangle$ with $v_1=(34)(56)$, $v_2=(12)(56)$, a Klein
four group. Let $W=\Ind_{H_4}^G\lambda_W$ with $\lambda_W(v_2)=-1$, and
$P=\Ind_{H_4}^G\triv$.
\item $H_8=\langle r,s\rangle$ with $r=(1234)(56)$, $s=(24)(56)$, dihedral of
order $8$. Let $Y=\Ind_{H_8}^G\lambda_Y$ with $\lambda_Y(r)=-1$,
$Z=\Ind_{H_8}^G\lambda_Z$ with $\lambda_Z(s)=-1$, and $C=\Ind_{H_8}^G\triv$.
\item $H_{10}=\langle c,d\rangle$ with $c=(12345)$, $d=(25)(34)$, dihedral of
order $10$. Let $V=\Ind_{H_{10}}^G\lambda_V$ with $\lambda_V(d)=-1$, and
$B=\Ind_{H_{10}}^G\triv$.
\item $H_{36}=\langle a,b,t\rangle$ with $a=(123)$, $b=(456)$,
$t=(14)(2536)$; it is $(C_3\times C_3)\rtimes C_4$, with $t^2$ acting by
inversion on $\langle a,b\rangle$. Let $U=\Ind_{H_{36}}^G\lambda_U$ with
$\lambda_U(t)=i$ and $A=\Ind_{H_{36}}^G\lambda_A$ with $\lambda_A(t)=-1$, both
trivial on $\langle a,b\rangle$.
\end{itemize}
All nine are genuine characters induced from explicitly specified linear
characters.

\begin{lemma}[The nine rows; \stU]\label[lemma]{grp:lem:a6-rows}
The following values hold on the six patterns; the last column is the
multiplicity of the trivial character.
\begin{center}
\begin{tabular}{@{}lrrrrrrr@{}}
\toprule
 & $1^6$ & $1^3\,3$ & $1^2\,2^2$ & $1\,5$ & $2\,4$ & $3^2$ & mean\\\midrule
$W$&90&0&$-2$&0&0&0&0\\
$Y$&45&0&1&0&$-1$&0&0\\
$P$&90&0&6&0&0&0&1\\\midrule
$U$&10&1&$-2$&0&0&1&0\\
$A$&10&1&2&0&$-2$&1&0\\
$V$&36&0&$-4$&1&0&0&0\\
$B$&36&0&4&1&0&0&1\\
$Z$&45&0&$-3$&0&1&0&0\\
$C$&45&0&5&0&1&0&1\\\bottomrule
\end{tabular}
\end{center}
\end{lemma}

\begin{proof}
Use the induction formula of \cref{grp:sec:conv}. In $H_4$ the three
nonidentity elements are double transpositions; a nontrivial linear character
sums to $-1$ on them and the trivial character to $3$. In $H_8$, $r^2$ and the
four reflections are double transpositions, and $r,r^3$ have type $2\,4$; the
character sums on these two types are $(1,-2)$ for $\lambda_Y$, $(-3,2)$ for
$\lambda_Z$ and $(5,2)$ for $\triv$. In $H_{10}$ the four nonidentity
rotations are $5$-cycles, two in each of the two classes, and the five
reflections are double transpositions. In $H_{36}$, $\langle a,b\rangle$ has
four single and four double $3$-cycles; the nine elements of
$\langle a,b\rangle t^2$ are involutions and the eighteen elements of the
cosets of $t$ and $t^3$ have order four, hence type $2\,4$. The contributions
of the last two cosets cancel for $\lambda_U$ and sum to $-18$ for
$\lambda_A$. Substitution gives every entry, and averaging over the class
sizes gives the last column.
\end{proof}

\begin{theorem}[Three positive character identities; \stU]\label{grp:thm:a6-identities}
With $\Reg$ the regular character of $A_6$,
\begin{align}
 5W+\triv&=\Reg+U+V+Z,\label{grp:eq:a6-W}\\
 12Y+4\,\triv&=\Reg+4A+V+3B,\label{grp:eq:a6-Y}\\
 10P+3\,\triv&=\Reg+3A+3B+9C.\label{grp:eq:a6-P}
\end{align}
\end{theorem}

\begin{proof}
Each equality holds entry by entry in \cref{grp:lem:a6-rows}, the regular
character being $360$ at the identity and zero elsewhere. An independent
verification enumerates the $360$ even permutations, constructs every
subgroup and linear character by closure and evaluates the induction formula
on every element in integer arithmetic; it is recorded in\\
\artifact{grh_a6_transfer_2026-11-25/character_certificate.json}.
\end{proof}

\begin{proposition}[Optimality in declared classes; \stU]\label[proposition]{grp:prop:a6-opt}
\begin{enumerate}[label=\textup{(\alph*)}]
\item $Y$ is the sum of the six nontrivial irreducible characters of $A_6$,
each with multiplicity one. Thus $45$ is the least degree of a genuine
character containing every nontrivial irreducible constituent.
\item Among integer-valued genuine characters of mean zero that are supported
on at most two of the six cycle patterns and contain every nontrivial
irreducible, the minimum degree is $90$, attained by $W$.
\item The coefficient $5$ in the order domination deduced from
\eqref{grp:eq:a6-W} is sharp at the level of all order functionals that are
nonnegative on genuine characters and vanish on $\triv$.
\end{enumerate}
\end{proposition}

\begin{proof}
(a) Take the inner product of \eqref{grp:eq:a6-Y} with a nontrivial
irreducible $\eta$: the right side has coefficient at least $\eta(1)>0$, so
$\langle Y,\eta\rangle\ge1$. The trivial coefficient of $Y$ is zero and
$\langle Y,Y\rangle=(45^2+45+90)/360=6$. There are six nontrivial irreducibles
because there are seven classes, so all six multiplicities equal one.

(b) The same argument with \eqref{grp:eq:a6-W} shows that $W$ contains every
nontrivial irreducible. Let $\xi$ be a genuine character of mean zero
containing all six, so $D=\xi(1)\ge45$ by (a). If its support contains only
one other pattern, of size $M$, its value there is the negative integer
$-D/M$; hence $M\mid D$ and
$\langle\xi,\xi\rangle=(D^2+D^2/M)/360\in\Z$, with
$M\in\set{40,45,144,90}$. For $45\le D<90$ the only possibilities are
$(D,M)=(80,40)$ and $(45,45)$, with squared norms $164/9$ and $23/4$, not
integers. A nonzero character of mean zero supported on one pattern does not
exist.

(c) The degree-five character $\sigma$, the natural permutation character
minus $\triv$, is irreducible by $2$-transitivity, and
$\sigma(1^2\,2^2)=1$. Hence
$\langle W,\sigma\rangle=(90\cdot5-2\cdot45)/360=1$ and
$\langle\Reg,\sigma\rangle=5$. The functional
$\xi\mapsto\langle\xi,\sigma\rangle$ is nonnegative on genuine characters,
vanishes on $\triv$ and attains the ratio $5$.
\end{proof}

Statement (c) does not assert that the extreme functional is realized by
actual Artin zeros, and the proposition does not rule out virtual characters,
several test functions at once, or other kernels. It shows that there is a
genuine tradeoff: the degree-$45$ detector $Y$ needs three patterns, and a
two-pattern detector of mean zero needs degree $90$.

\subsubsection{Zero detectors and the prime criterion}
Let $L/\Q$ be Galois with group $A_6$. The character $W$ corresponds to a
quadratic character of $L^{\ker\lambda_W}/L^{H_4}$, of degrees $180$ and $90$
over $\Q$, so $L(s,W)=\zeta_{L^{\ker\lambda_W}}(s)/\zeta_{L^{H_4}}(s)$ is
entire; likewise $Y$ is the quadratic relative factor of degrees $90$ and
$45$, and $L(s,P)=\zeta_{L^{H_4}}(s)$. The character $U$ comes from a cyclic
quartic character over the degree-ten field $L^{H_{36}}$; it is real-valued
and irreducible, since $\langle U,U\rangle=(100+40+180+40)/360=1$.

\begin{theorem}[Multiplicity domination; \stU]\label{grp:thm:a6-orders}
For $\rho\ne1$ write $m_J=\ord_\rho L(s,J)$, $m_L=\ord_\rho\zeta_L$ and
$m_0=\ord_\rho\zeta$. Then
\[
 0\le m_L\le5m_W+m_0,\qquad
 0\le m_L\le12m_Y+4m_0,\qquad
 0\le m_L\le10m_P+3m_0 .
\]
Given the Riemann hypothesis for $\zeta$, in $1/2<\Re\rho<1$ one has $m_0=0$,
and for each $J\in\set{W,Y,P}$ the following are equivalent: the Riemann
hypothesis for $\zeta_L$; the absence of zeros of $L(s,J)$ in $1/2<\Re s<1$.
\end{theorem}

\begin{proof}
By Artin formalism \eqref{grp:eq:a6-W}--\eqref{grp:eq:a6-P} are identities of
meromorphic functions, for example
$L(s,W)^5\zeta(s)=\zeta_L(s)L(s,U)L(s,V)L(s,Z)$. Take orders and drop the
nonnegative orders of the auxiliary Hecke and Dedekind functions, none of which
has a pole at $\rho\ne1$. This gives the inequalities and the implication from
$L(s,J)$ to $\zeta_L$. The converse is \eqref{grp:eq:heil}, each $J$ being a
genuine character.
\end{proof}

\begin{theorem}[Single-detector prime criterion; \stR]\label{grp:thm:a6-criterion}
Use the Riemann hypothesis for $\zeta$ from \cref{src:thm:dgrh}. For each
$J\in\set{W,Y,P}$ separately, with $\delta_W=\delta_Y=0$ and $\delta_P=1$, the
following are equivalent.
\begin{enumerate}[label=\textup{(\roman*)}]
\item $\zeta_L$ satisfies the Riemann hypothesis.
\item For every $\eps>0$, $R_J(x)=O_{L,\eps}(x^{-1/2+\eps})$.
\item For every $\eps>0$ there is $K_\eps$ with
$R_J(x)\le K_\eps x^{-1/2+\eps}$ for all sufficiently large real $x$.
\item For every $\eps>0$ there is $K_\eps$ with
$P_J(n)-P_J(2n)+\delta_J\log2\le K_\eps n^{-1/2+\eps}$ for all sufficiently
large integers $n$.
\end{enumerate}
The excluded set $S$ may be any fixed finite set containing the ramified
primes. If these conditions hold, every Artin $L$-function of $L/\Q$ is
holomorphic and nonzero in $1/2<\Re s<1$.
\end{theorem}

\begin{proof}
Each $J$ is monomially certified, and \cref{grp:thm:a6-orders} supplies the
domination required in \cref{grp:cor:detector}(ii), with $k=5,12,10$.
\end{proof}

\begin{example}[A specified sextic]\label[example]{grp:ex:a6-field}
Let
\begin{equation}\label{grp:eq:a6-f}
 f(X)=X^6-2X^5-X^4+2X^2-1 ,
\end{equation}
listed as an $A_6$ polynomial in \cite{DokchitserGroupNames}. Its discriminant
is $287296=2^6\cdot67^2$, a square. Its reductions factor as
\[
\begin{array}{r|l}
3&(X+1)(X^5-X^3+X^2+X-1)\\
5&(X^2-2X-2)(X^4+X^2+2X-2)\\
109&(X+20)(X-48)(X-13)(X^3+39X^2+14X+2)
\end{array}
\]
with irreducible factors
(\artifact{grh_a6_transfer_2026-11-25/field_check.txt}). A proper rational
factor would have degree $1$ or $5$ by the reduction at $3$ and degree $2$ or
$4$ by the reduction at $5$; so $f$ is irreducible. The Galois group is
transitive, lies in $A_6$, and contains a $5$-cycle. A transitive subgroup of
$S_6$ containing a $5$-cycle is primitive: a nontrivial block system would
have $2$ or $3$ blocks, the $5$-cycle would act trivially on the set of blocks
and then on each block. The factorization at $109$ gives a $3$-cycle, and
Jordan's theorem on primitive groups containing a cycle of prime length
$p\le n-3$ gives $A_6\subseteq G$. Hence $G=A_6$.

Take $S=\set{2,67}$ and classify $p\notin S$ by the degrees of the irreducible
factors of $f$ modulo $p$. Each of the following integral weights, used alone
in \cref{grp:thm:a6-criterion}(iv) with the indicated $\delta$, is equivalent
to the Riemann hypothesis for the splitting field of $f$, of degree $360$:
\begin{center}
\begin{tabular}{@{}lrrrr@{}}
\toprule
 & $1^6$ & $1^2\,2^2$ & $2\,4$ & $\delta$\\\midrule
two patterns, signed ($W/2$)&45&$-1$&0&0\\
three patterns, signed ($Y$)&45&1&$-1$&0\\
two patterns, positive ($P/6$)&15&1&0&$1/6$\\\bottomrule
\end{tabular}
\end{center}
All other patterns have weight zero. In the first row the criterion reads
\[
 S(n)-S(2n)\le C_\eps n^{-1/2+\eps},\qquad
 S(x)=\sum_{p\le x}\bigl(45\,\grpone_{1^6}(p)-\grpone_{1^2 2^2}(p)\bigr)
 \log p\Bigl(\frac1p-\frac1x\Bigr),
\]
with no centring constant and two prime patterns. Among the primes
$p\le5000$, $p\notin S$, the six patterns occur as follows
(\artifact{grh_a6_transfer_2026-11-25/sample_patterns.csv}):
\begin{center}
\begin{tabular}{@{}lrrrrrr@{}}
\toprule
pattern & $1^6$ & $1^3\,3$ & $1^2\,2^2$ & $1\,5$ & $2\,4$ & $3^2$\\\midrule
count & 3 & 83 & 74 & 262 & 174 & 71\\
group density & $1/360$ & $1/9$ & $1/8$ & $2/5$ & $1/4$ & $1/9$\\\bottomrule
\end{tabular}
\end{center}
The small number of completely split primes shows why numerical
experimentation cannot replace a proof of the estimate.
\end{example}

\begin{remark}[The involution filter]\label[remark]{grp:rem:a6-filter}
For $F=L^{H_4}$ and the quadratic extension $L^{\ker\lambda_W}/F$, the value
$P(\Frob_p)$ counts the degree-one primes of $F$ above an unramified $p$, and
$W(\Frob_p)$ sums the quadratic character over them. The nonnegative
degree-one filter analogous to the source's $1-\lambda$ is
\[
 P-W=8\,\grpone_{1^2 2^2},
\]
which selects the involution class and nothing else. It is, however, a
virtual character and not an entire Hecke factor:
$\langle P-W,\sigma\rangle=1$ while $\langle P-W,U\rangle=-2$, so even the cone
of genuine order characters permits both signs for this count, and the
positivity of real residues required in \cref{grp:prop:interface}(b) is not
available for it. In contrast $P$ has positive coefficients and the correct
order sign, and $W$ has signed coefficients and the correct order sign.
\end{remark}

As for $A_5$, the group $A_6$ is perfect, so by \cref{nab:thm:periodic} no function of
$p$ modulo a fixed modulus equals a nonconstant weight of the table, even in mean square.

\subsection{The simple group of order \texorpdfstring{$168$}{168}}\label{grp:sec:psl27}
{\rightskip=0pt plus 1fil\relax
\sources{grh_sylow_transfer_2026-11-25}
}

Let $G=\GL_3(\F_2)\simeq\PSL_2(\F_7)$. Its classes have element orders
$1,2,3,4,7,7$, and its character table \cite{DokchitserGroupNames} is
\begin{center}
\begin{tabular}{crrrrrr}
\toprule
 & $1$ & $2$ & $3$ & $4$ & $7A$ & $7B$\\
size & 1 & 21 & 56 & 42 & 24 & 24\\\midrule
$\triv$ & 1 & 1 & 1 & 1 & 1 & 1\\
$\chi_3$ & 3 & $-1$ & 0 & 1 & $\alpha$ & $\bar\alpha$\\
$\overline\chi_3$ & 3 & $-1$ & 0 & 1 & $\bar\alpha$ & $\alpha$\\
$\chi_6$ & 6 & 2 & 0 & 0 & $-1$ & $-1$\\
$\chi_7$ & 7 & $-1$ & 1 & $-1$ & 0 & 0\\
$\chi_8$ & 8 & 0 & $-1$ & 0 & 1 & 1\\
\bottomrule
\end{tabular}
\qquad $\alpha=(-1+\sqrt{-7})/2$.
\end{center}
Let $P=\langle r,s:r^4=s^2=1,\ srs=r^{-1}\rangle$ be a Sylow $2$-subgroup, with
quadratic characters $\lambda_A(r)=1$, $\lambda_A(s)=-1$ and
$\lambda_B(r)=-1$, $\lambda_B(s)=1$. Let $J\simeq S_4$ be the stabilizer of a
nonzero vector of $\F_2^3$ and $U\simeq C_7\rtimes C_3$ the normalizer of a
Sylow $7$-subgroup, with a nontrivial character $\omega$ of $U/C_7$.

\begin{lemma}[Induction identities; \stU]\label[lemma]{grp:lem:psl27-ind}
\begin{align}
 \Ind_P^G\lambda_A&=\chi_3+\overline\chi_3+\chi_7+\chi_8,&
 \Ind_P^G\lambda_B&=\chi_6+\chi_7+\chi_8,\label{grp:eq:psl27-A}\\
 \Ind_J^G\triv&=\triv+\chi_6,\qquad \Ind_J^G\grpsgn=\chi_7,&
 \Ind_U^G\omega&=\chi_8.\label{grp:eq:psl27-B}
\end{align}
Hence $W_P=\Ind_P^G\lambda_A+\Ind_P^G\lambda_B$ has the values
$(42,-2,0,0,0,0)$ in table order, and $2W_P$ is induced from the nontrivial
character of a subgroup of order two.
\end{lemma}

\begin{proof}
Write $c_j$ for the value of a character of $G$ on the class of order $j$. The
subgroup $P$ contains five involutions and two elements of order four, on
which the two linear characters have sums $(-3,2)$ and $(1,-2)$; so the inner
products with the two inductions are $(c_1-3c_2+2c_4)/8$ and
$(c_1+c_2-2c_4)/8$. For $J$ the cycle counts of $S_4$ give
$(c_1+9c_2+8c_3+6c_4)/24$ and, for the sign character,
$(c_1-3c_2+8c_3-6c_4)/24$. For $U$ the fourteen elements outside $C_7$
contribute $7(\omega+\bar\omega)c_3=-7c_3$ and the six elements of order seven
contribute $3c_{7A}+3c_{7B}$; divide by $21$. Evaluating on the table gives
the multiplicities. The computation is verified in exact arithmetic on the
group of invertible binary matrices in
\artifact{grh_sylow_transfer_2026-11-25/group_certificate.json}.
\end{proof}

The product $F_P(s)=L(s,\Ind_P^G\lambda_A)L(s,\Ind_P^G\lambda_B)$ is a product
of two quadratic Hecke $L$-functions over the field $L^P$ of degree $21$. Its
prime weight $W_P/2$ is $21$ at a completely split prime, $-1$ at a prime
whose Frobenius is an involution, and zero elsewhere.

\begin{theorem}[Two-pattern test for $\PSL_2(7)$; \stR]\label{grp:thm:psl27}
Use the Riemann hypothesis for $\zeta$. Let $L/\Q$ be Galois with group
$\PSL_2(7)$ and $w=W_P/2$, so $\delta_w=0$. At every $\rho$ with
$1/2<\Re\rho<1$,
\[
 \ord_\rho\zeta_L\le6\,\ord_\rho F_P ,
\]
and the Riemann hypothesis for $\zeta_L$ is equivalent to
\begin{equation}\label{grp:eq:psl27-test}
 21\sum_{\substack{p\le x\\\Frob_p=1}}\log p\Bigl(\frac1p-\frac1x\Bigr)
 -\sum_{\substack{p\le x\\\ord(\Frob_p)=2}}\log p\Bigl(\frac1p-\frac1x\Bigr)
 -c_w\le C_{L,\eps}x^{-1/2+\eps}\qquad(x\ge x_\eps)
\end{equation}
for every $\eps>0$, the sums omitting a fixed finite set containing the
ramified primes. The test does not presuppose Artin holomorphy.
\end{theorem}

\begin{proof}
At such $\rho$ write $n_1,n_3,n_{\bar3},n_6,n_7,n_8$ for the Artin orders,
$u=n_3+n_{\bar3}$ and $d=\ord_\rho\zeta_L$. The Riemann hypothesis for $\zeta$
gives $n_1=0$. By \eqref{grp:eq:psl27-B}, $n_6=\ord_\rho\zeta_{L^J}\ge0$,
$n_7=\ord_\rho L(s,\grpsgn,L/L^J)\ge0$ and
$n_8=\ord_\rho L(s,\omega,L/L^U)\ge0$; nothing is assumed about the signs of
$n_3,n_{\bar3}$. The two Hecke orders in \eqref{grp:eq:psl27-A},
$a_P=u+n_7+n_8$ and $b_P=n_6+n_7+n_8$, are nonnegative. Taking degrees in the
regular character,
\[
 d=3u+6n_6+7n_7+8n_8=3a_P+6b_P-2n_7-n_8\le6(a_P+b_P).
\]
Thus $F_P$ and $\zeta_L$ have the same zeros in the right half of the strip,
the converse inclusion being \eqref{grp:eq:heil}. Now apply
\cref{grp:cor:detector}(ii) to $J=W_P$ and divide the inequality by $2$.
\end{proof}

\begin{example}[Trinks' polynomial]\label[example]{grp:ex:trinks}
The polynomial $g(X)=X^7-7X+3$ has Galois group $\PSL_2(7)$
\cite{ErbachFischerMcKay1979Polynomials} and discriminant $3^87^8$. In the
action on seven points the cycle types are $1^7$, $1^3\,2^2$, $1\,3^2$,
$1\,2\,4$ and $7$. In \eqref{grp:eq:psl27-test} the first sum is therefore
over the primes $p\ne3,7$ modulo which $g$ splits into linear factors, and
the second over those with three linear and two irreducible quadratic
factors; all other primes, including those with irreducible reduction, are
ignored. The factorization \eqref{grp:eq:psl27-A} reduces the degree of the
base field of the Hecke characters from $84$ to $21$ without changing the
test. Since $G$ is simple, the weight $W_P$ is not a function of $p$ modulo any fixed
modulus (\cref{nab:thm:periodic}).
\end{example}

The general Sylow-filtered race \cref{nab:thm:sylow} for this group uses all
classes of prime-power order; the two-pattern test is more selective and has
the explicit multiplicity constant $6$.

\subsection{\texorpdfstring{$\PSL_2(q)$}{PSL2(q)} for even \texorpdfstring{$q$}{q}}\label{grp:sec:psl-even}
{\rightskip=0pt plus 1fil\relax
\sources{psl2_grh_transfer}
}

Let $q=2^r\ge4$ and $G=\PSL_2(q)=\SL_2(q)=\PGL_2(q)$, of order $q(q^2-1)$,
acting on $X=\mathbb P^1(\Fq)$. Let $B$ be a point stabilizer, $U\lhd B$ the
group of translations, $T$ the split torus of order $q-1$ and $T'$ a nonsplit
torus of order $q+1$; their normalizers $N,N'$ are dihedral of orders
$2(q-1)$ and $2(q+1)$, and $B/U\simeq\Fq^\times$. Write $f(h)=\abs{X^h}$. The
conjugacy types follow from the characteristic polynomial: the nontrivial
semisimple classes are indexed by pairs of inverse eigenvalues in $\Fq^\times$
or in the norm-one subgroup of $\F_{q^2}^\times$, with the tori as
centralizers; the centralizer of a nonidentity unipotent is $U$, and all
nonidentity unipotents are conjugate, since squaring is bijective on
$\Fq^\times$. In particular there are $q+1$ classes.
\begin{center}\small
\begin{tabular}{@{}lrrrr@{}}\toprule
type & $f$ & centralizer & classes & elements\\\midrule
identity & $q+1$ & $q(q^2-1)$ & $1$ & $1$\\
unipotent $u\ne1$ & $1$ & $q$ & $1$ & $q^2-1$\\
split semisimple & $2$ & $q-1$ & $(q-2)/2$ & $q(q+1)(q-2)/2$\\
nonsplit semisimple & $0$ & $q+1$ & $q/2$ & $q^2(q-1)/2$\\\bottomrule
\end{tabular}
\end{center}
The projective-line character is $\pi_B=f=\triv+\grpSt$ with $\grpSt$
irreducible of degree $q$, the action being doubly transitive. Let $\lambda$
be a nontrivial character of $U$; its values are $\pm1$. Define
\[
 \Gamma=\Ind_U^G\lambda,\qquad W=\triv+\Gamma,
\]
let $P$ be the sum of $\Ind_B^G\theta$ over one representative of each inverse
pair of nontrivial characters $\theta$ of $B/U$, and let
$D=\Ind_{N'}^G\eta$, where $\eta$ is $1$ on $T'$ and $-1$ on
$N'\setminus T'$. These are genuine characters.

\begin{lemma}[Character identities; \stU]\label[lemma]{grp:lem:psl-identities}
\begin{align}
 \Gamma&=\grpSt+P+D, & W&=\triv+\grpSt+P+D,\label{grp:eq:psl-gamma}\\
 \pi_T&=\triv+2\grpSt+P+D, & \pi_{T'}&=\triv+P+D,\label{grp:eq:psl-perms}\\
 \Reg&=\triv+q\,\grpSt+(q+1)P+(q-1)D.\label{grp:eq:psl-reg}
\end{align}
In particular $W=\pi_B+\pi_{T'}-\triv$ and $\pi_T=W+\grpSt$. The values on the
four types are
\begin{center}\small
\begin{tabular}{@{}lrrrr@{}}\toprule
 & identity & unipotent & split & nonsplit\\\midrule
$\grpSt$ & $q$ & $0$ & $1$ & $-1$\\
$P$ & $(q-2)(q+1)/2$ & $(q-2)/2$ & $-1$ & $0$\\
$D$ & $q(q-1)/2$ & $-q/2$ & $0$ & $1$\\
$\Gamma$ & $q^2-1$ & $-1$ & $0$ & $0$\\
$W$ & $q^2$ & $0$ & $1$ & $1$\\
$\pi_T$ & $q(q+1)$ & $0$ & $2$ & $0$\\
$\pi_{T'}$ & $q(q-1)$ & $0$ & $0$ & $2$\\\bottomrule
\end{tabular}
\end{center}
\end{lemma}

\begin{proof}
Use the induction formula. For $\Gamma$ the unipotent sum in $U$ is $-1$
because $\lambda$ is nontrivial, and no nonidentity semisimple element lies in
$U$. Induction from $B$ is computed by summing over fixed points: a split
element with eigenvalue ratio $a\ne1$ contributes $\theta(a)+\theta(a^{-1})$
and a unipotent contributes $1$. Thus the induction is unchanged under
$\theta\mapsto\theta^{-1}$, and $P=(\pi_U-\pi_B)/2$ by induction in stages
from the regular character of $B/U$. For $D$, $N'$ has $q+1$ involutions, all
of sign $-1$, and a nonsplit semisimple class meets $T'$ in two elements of
sign $1$. For $\pi_T$ count ordered pairs of distinct fixed points; for
$\pi_{T'}$ use the same two-element intersection. The table gives each
identity by substitution.
\end{proof}

\begin{theorem}[Optimal nonnegative detector; \stU]\label{grp:thm:psl-detector}
$\Gamma$ is the sum of all nontrivial irreducible characters of $G$, each with
multiplicity one. Among genuine characters that contain every nontrivial
irreducible and take only nonnegative values, $W$ has the least degree, $q^2$,
and is the unique minimizer.
\end{theorem}

\begin{proof}
$G$ is simple for $q\ge4$. In a nontrivial irreducible representation $U$
cannot act trivially, since its normal closure is $G$; so the restriction to
$U$ contains a nontrivial additive character. Conjugation by $T$ permutes the
nontrivial additive characters transitively, so the chosen $\lambda$ occurs.
By Frobenius reciprocity all $q$ nontrivial irreducibles occur in $\Gamma$ and
the trivial one does not. Its squared norm is
\[
 \frac{(q^2-1)^2+(q^2-1)}{q(q^2-1)}=q,
\]
and a sum of $q$ positive integer squares equals $q$ only if each is one.
Hence the least degree of a character containing the nontrivial irreducibles
is $\Gamma(1)=q^2-1$, with equality only for $\Gamma$, whose value at $u$ is
$-1$. Adding the trivial character gives $W\ge0$; at degree $q^2$ the extra
dimension is one, and the simple nonabelian group $G$ has no other
one-dimensional character.
\end{proof}

The argument does not use the tabulated character table of $\PSL_2(q)$. It
shows that $\grpSt$, $P$, $D$ partition the nontrivial irreducibles, with
multiplicity one; by \eqref{grp:eq:psl-reg} their constituents have degrees
$q$, $q+1$, $q-1$, so $P$ contains $(q-2)/2$ irreducibles and $D$ contains
$q/2$.

\begin{theorem}[Exact transfers; \stG]\label{grp:thm:psl-transfer}
Let $L/\Q$ be Galois with group $G$ and $F_H=L^H$, and use the Riemann
hypothesis for $\zeta$ from \cref{src:thm:dgrh}. The following are
equivalent:
\begin{enumerate}[label=\textup{(\arabic*)}]
\item the Riemann hypothesis for $\zeta_L$;
\item the Riemann hypothesis for $\zeta_{F_T}$, of degree $q(q+1)$;
\item the Riemann hypothesis for both $\zeta_{F_B}$ and $\zeta_{F_{T'}}$, of
degrees $q+1$ and $q(q-1)$;
\item $L(s,W)$ has no zeros in $1/2<\Re s<1$;
\item the Hecke $L$-function $L(s,\lambda,L/F_U)$ has no zeros in
$1/2<\Re s<1$.
\end{enumerate}
More precisely, at any $\rho$ with $1/2<\Re\rho<1$,
\begin{align}
 \frac q2\ord_\rho\zeta_{F_T}&\le\ord_\rho\zeta_L\le(q+1)\ord_\rho\zeta_{F_T},
 \label{grp:eq:psl-ineq1}\\
 (q-1)\ord_\rho L(s,W)&\le\ord_\rho\zeta_L\le(q+1)\ord_\rho L(s,W),
 \label{grp:eq:psl-ineq2}
\end{align}
and
\[
 L(s,W)=\zeta(s)L(s,\lambda,L/F_U)
 =\frac{\zeta_{F_B}(s)\zeta_{F_{T'}}(s)}{\zeta(s)} .
\]
If these conditions hold, every Artin $L$-function of $L/\Q$ is holomorphic
and nonzero in $1/2<\Re s<1$.
\end{theorem}

\begin{proof}
Let $s_0=\ord_\rho L(s,\grpSt)$, $p_0=\ord_\rho L(s,P)$,
$d_0=\ord_\rho L(s,D)$. At $\rho$, $\zeta$ is finite and nonzero, so
$s_0=\ord_\rho\zeta_{F_B}\ge0$. The character $P$ is a sum of characters
induced from nontrivial linear characters and $D$ is induced from the
nontrivial character $\eta$, so $p_0,d_0\ge0$; nothing about individual Artin
factors is used. By \eqref{grp:eq:psl-gamma}--\eqref{grp:eq:psl-reg},
\begin{align*}
 \ord_\rho\zeta_L&=qs_0+(q+1)p_0+(q-1)d_0, &
 \ord_\rho\zeta_{F_T}&=2s_0+p_0+d_0,\\
 \ord_\rho\zeta_{F_{T'}}&=p_0+d_0,\qquad\ord_\rho\zeta_{F_B}=s_0, &
 \ord_\rho L(s,W)&=s_0+p_0+d_0=\ord_\rho L(s,\lambda,L/F_U).
\end{align*}
All equivalences and inequalities follow from the nonnegativity of the three
integers and $q\ge4$, together with the functional equations of the Dedekind
zeta functions and their nonvanishing on $\Re s\ge1$. The displayed identities
for $L(s,W)$ are \eqref{grp:eq:psl-gamma} and $W=\pi_B+\pi_{T'}-\triv$; the
first shows holomorphy away from a simple pole at $1$, which the nonvanishing
of the Hecke function at $1$ does not cancel. The last assertion is
\eqref{grp:eq:heil}.
\end{proof}

A family of subfields $F_{H_i}$ of $L$ is called \emph{spectrally detecting}
if the sum of its permutation characters contains every nontrivial irreducible
character. The condition tests which Artin components are visible; it does not
by itself exclude cancellation between visible meromorphic factors.
\Cref{grp:thm:psl-transfer} supplies the analytic noncancellation for the
optimal families of the next theorem.

\begin{theorem}[Sharp degrees; \stU]\label{grp:thm:psl-optimal}
For every $q=2^r\ge4$:
\begin{enumerate}[label=\textup{(\alph*)}]
\item the least degree of a single spectrally detecting subfield is $q(q+1)$;
\item the least possible maximum degree in a spectrally detecting family of
subfields is $q(q-1)$, attained by two subfields of degrees $q+1$ and
$q(q-1)$.
\end{enumerate}
\end{theorem}

\begin{proof}
The upper bounds follow from \eqref{grp:eq:psl-perms} and
\cref{grp:thm:psl-detector}. For the lower bounds we show that a subgroup $H$
of order greater than $q-1$ misses at least one nontrivial irreducible, and a
subgroup of order greater than $q+1$ misses all of $D$, meaning
$\langle D,\pi_H\rangle=\dim D^H=0$. We use Dickson's classification in the
form of \cite[Theorem~A.1]{GuralnickZieve2010Polynomials}; only the list of
isomorphism types and the element types of their embeddings are needed.

\emph{Subgroups of Borel type.} If $H$ has a nontrivial normal elementary
abelian $2$-subgroup $V$, it fixes a projective point: every involution has a
unique fixed point, commuting involutions share it, and normality makes $H$
preserve it. Conjugate into $B$. If $H\cap U$ is trivial, $H$ is cyclic of
order dividing $q-1$. Otherwise put $V=H\cap U$ and let the image of $H$ in
$B/U$ have order $d$. If $V\ne U$, then $V$ is a proper vector subspace over
the subfield $\F_{q_0}$ generated by the image scalars, so $\abs V\le q/q_0$,
$d\le q_0-1$ and $\abs H=\abs Vd\le q-q/q_0\le q-1$. If $V=U$, then $D^H=0$
because the sum of $D$ over $U$ is $q(q-1)/2-(q-1)q/2=0$.

\emph{Cyclic and dihedral types.} A cyclic subgroup of odd order has order
dividing $q-1$ or $q+1$, and a dihedral subgroup has order $2d$ with
$d\mid q\pm1$. A proper divisor $d$ is at most $(q\pm1)/3$, so $2d\le q-1$ for
$q\ge8$, and the same is checked directly for $q=4$. The only cases larger
than $q-1$ are the full nonsplit torus and the two full dihedral groups. On a
nonsplit torus the sum of $\grpSt$ is $q-q=0$. On the two full dihedral groups
the sums of $D$ are
\[
 \frac{q(q-1)}2-(q-1)\frac q2=0,\qquad
 \frac{q(q-1)}2-(q+1)\frac q2+q=0 .
\]
The nonsplit torus has order exactly $q+1$, so it does not affect (b).

\emph{Subfield and exceptional types.} A subgroup isomorphic to
$\PSL_2(q_0)$ or $\PGL_2(q_0)$ with $q=q_0^k$, $k\ge2$, has order
$q_0(q_0^2-1)$. If $k\ge3$ this is at most $q-q_0<q-1$. If $k=2$ it has
$q_0^2-1=q-1$ involutions, and every other nonidentity element has odd order
dividing $q_0\pm1$, hence dividing $q-1$, and is split in $G$; the sum of $D$
is again zero. The group $A_4$ has a normal Klein four group and lies in a
Borel subgroup. $S_4$ does not occur, as $G$ has no element of order four.
$A_5$ requires $5\mid q^2-1$, so $r$ is even: for $q=4$ it is $G$, for $q=16$
it is the square-subfield type just treated, and for $q\ge64$ its order $60$
is at most $q-1$. Finally $G$ itself has no nontrivial fixed constituents.

Since $[F_H:\Q]=\abs G/\abs H$, degrees less than $q(q+1)$ correspond to
$\abs H>q-1$, and in a family of degrees less than $q(q-1)$ every subgroup has
order greater than $q+1$ and misses $D$.
\end{proof}

The statement concerns subfields of the given $L$, not auxiliary extensions or
other analytic identities. The spectral obstruction is also an obstruction to
a purely formal nonnegative-order argument: assign order one to an invisible
irreducible and zero to the others; the selected permutation characters have
order zero while the regular character has positive order. This assignment is
not claimed to be realized by Artin zeros. The finite checks for $q=4,8$ (all
$59$, respectively $386$, subgroups) and of the identities for $q=16$ are in
\artifact{psl2_grh_transfer/checks.json}.

\begin{corollary}[Root-count criterion; \stR]\label[corollary]{grp:cor:psl-root}
Use the Riemann hypothesis for $\zeta$. Let $F(X)$ be a monic integral minimal
polynomial of a primitive integral element of $F_B$; it has degree $q+1$ and
splitting field $L$. For $p\nmid\disc F$ let $f_p$ be the number of roots of
$F$ modulo $p$, and define the weight $w_p$ by
\[
\begin{array}{c|cccc}
 f_p&q+1&1&2&0\\\hline
 w_p&q^2&0&1&1
\end{array}
\]
Then $w_p=W(\Frob_p)$, the four cases exhaust the good primes, and the Riemann
hypothesis for $\zeta_L$ is equivalent to
\[
 R_W(x)=\sum_{p\le x,\ p\nmid\disc F}w_p\log p\Bigl(\frac1p-\frac1x\Bigr)
 -\log x-c_W\le C_\eps x^{-1/2+\eps}\qquad(x\ge x_\eps)
\]
for every $\eps>0$. Under the Riemann hypothesis for $\zeta_L$ one has
$R_W(x)=O_L(x^{-1/2}\log^2(2x))$. The weights are nonnegative, and
$0\le a_W(p^k)\le q^2\log p$ at all prime powers.
\end{corollary}

\begin{proof}
$B$ has trivial core, so the splitting field of $F$ is $L$, and
$f_p=f(\Frob_p)$ by the Frobenius factorization theorem. The character $W$ is
monomially certified with $a=1$, and \eqref{grp:eq:psl-ineq2} supplies the
domination in \cref{grp:cor:detector}(ii) with $k=q+1$.
\end{proof}

Subtracting the ordinary prime sum, which is controlled by the Riemann
hypothesis for $\zeta$, gives the signed contrast with weight $q^2-1$ for
complete splitting, $-1$ for exactly one root and $0$ otherwise. This weight
is $\Gamma$; it has mean zero, and the logarithmic main term disappears. The
two formulations are equivalent; the positive one is the natural target for a
construction based on divisibility. By \cref{nab:thm:periodic}, $G$ being perfect, the
least limiting mean-square error of predicting $w_p$ by a function of $p$ modulo a fixed
modulus is exactly $\norm{\Gamma}^2=q$, and the optimal periodic prediction is the
constant $1$. The extension of \cref{grp:thm:psl-transfer,grp:cor:psl-root} to the
relative extension $LK/K$ over an abelian field $K$ is \cref{nab:rem:basechange}.

\begin{example}\label[example]{grp:ex:psl-table}
The first cases are
\begin{center}
\begin{tabular}{@{}rrrrrr@{}}\toprule
$q$ & $[L:\Q]$ & single subfield & pair of subfields & $W(1)$ & $\Gamma(1)$\\\midrule
4 & 60 & 20 & 5, 12 & 16 & 15\\
8 & 504 & 72 & 9, 56 & 64 & 63\\
16 & 4080 & 272 & 17, 240 & 256 & 255\\
32 & 32736 & 1056 & 33, 992 & 1024 & 1023\\\bottomrule
\end{tabular}
\end{center}
For $q=4$ this is the detector $W_{16}$ of \cref{grp:sec:a5}. For $q=8$ any
polynomial of degree nine with Galois group $\PSL_2(8)$ in its natural action
(transitive group 9T27, see \cite{KlunersMalleDatabase}) is eligible: the four
types have $1,63,216,224$ elements, the weights $64,0,1,1$ have mean one, and
the problem for the field of degree $504$ becomes one root-count inequality at
rational primes. No particular polynomial is selected or certified here.
\end{example}

\subsection{\texorpdfstring{$\PSL_2(q)$}{PSL2(q)} for odd \texorpdfstring{$q$}{q}}\label{grp:sec:psl-odd}
{\rightskip=0pt plus 1fil\relax
\sources{grh_transfers_density_gaps_and_zero_tests; psl2_grh_transfer; ProveIt_research_2026-11-27}
}

The transitivity of the torus on nontrivial additive characters, used in
\cref{grp:thm:psl-detector}, fails in odd characteristic. The least dimension
of a nonnegative detector and the least degree of a single spectrally
detecting subfield can nevertheless be determined. Let $q\ge5$ be an odd prime
power, $G=\PSL_2(\Fq)$, and let $T_s,T_n$ be the split and nonsplit projective
tori, of orders $a=(q-1)/2$ and $b=(q+1)/2$. Put $D=q(q+1)/2$, so that
$\abs G=(q-1)D$, and let
\[
 S=\sum_{\chi\in\grpIrr(G)}\chi
\]
be the Gelfand character. We use the ordinary character table of
$\SL_2(q)$ \cite[Table~1]{BohmlerFarrellLassueur2022Trivial}, restricted to
the rows trivial on $-I$. Besides $\triv$ and the Steinberg character of
degree $q$, with values $0,1,-1$ on nontrivial unipotents, split and nonsplit
elements, there are: the principal series, indexed by nontrivial characters
$\alpha$ of $T_s$ up to inversion, an inverse pair contributing one row of
degree $q+1$ with values $1$, $\alpha(t)+\alpha(t)^{-1}$, $0$ on the three
nonidentity types; and the cuspidal series, indexed by nontrivial characters
$\beta$ of $T_n$ up to inversion, a pair contributing one row of degree $q-1$
with values $-1$, $0$, $-\beta(t)-\beta(t)^{-1}$. A quadratic parameter
contributes two exceptional rows of half the corresponding degree, whose sum
has the displayed series values for that parameter.

\begin{theorem}[Least nonnegative dimension; \stU]\label{grp:thm:odd-gelfand}
Let $\delta$ be the quadratic character of whichever projective torus has even
order. The values of $S$ are
\[
\begin{array}{c|cccc}
q\bmod4 &\text{identity}&\text{nontrivial unipotent}
 &\text{split }t\ne1&\text{nonsplit }t\ne1\\ \hline
1&D+1&1&1+\delta(t)&1\\
3&D&0&1&1-\delta(t)
\end{array}
\]
toral involutions included. In particular $S\ge0$, and $S$ is the unique
genuine character of least degree among those that take only nonnegative
values and contain every nontrivial irreducible character.
\end{theorem}

\begin{proof}
If $q=4k+1$ the numbers of ordinary inverse pairs in the two series are $k-1$
and $k$, and the exceptional pair is principal; if $q=4k+3$ they are $k$ and
$k$, and the exceptional pair is cuspidal. Summing the degrees and the
unipotent values gives the first two columns. For nontrivial $t$ the sum of
all nontrivial characters of its cyclic torus is $-1$; a quadratic parameter,
when present, contributes twice its value to the parameter sum in the table,
hence one extra copy beyond this $-1$. Adding the trivial and Steinberg terms
gives the remaining columns. A nonzero pointwise nonnegative genuine character
has strictly positive mean, its value at the identity being its dimension, so
it contains $\triv$. Covering all nontrivial irreducibles then forces degree
at least $S(1)$, with equality exactly for $S$.
\end{proof}

For $q=5,7,11,13,17,19$ the minima are $16,28,66,92,154,190$; the case $q=5$
is again $W_{16}$. The values of $S$ are verified for these six fields by
exact arithmetic on projective matrices, independently of the character
table; the output is\\
\artifact{grh_transfers_density_gaps_and_zero_tests/verification/check_odd_psl2.txt}.\\
The verification writes $S=\sum m_\chi\chi$ for the proposed class function;
by orthogonality $(S*S)(g)=\abs G\sum m_\chi^2\chi(g)/\chi(1)$ for the
unnormalized convolution, while the commutator-count function has the same
expansion with every $m_\chi^2$ replaced by one. The checker counts
commutators, verifies the equality of the two on every class, and verifies
positive definiteness of convolution by $S$ on class functions, whose
eigenvalues are $\abs Gm_\chi/\chi(1)$; hence every $m_\chi=1$.

\begin{proposition}[A Hecke realization for $q\equiv1\pmod4$; \stU]
\label[proposition]{grp:prop:odd-hecke}
Let $q\equiv1\pmod4$, let $N=T_s\rtimes\langle w\rangle$ be the dihedral
normalizer of $T_s$, and extend $\delta$ to a character $\lambda$ of $N$ by
$\lambda(w)=1$. Then
\[
 S=\triv+\Ind_N^G\lambda .
\]
Consequently, for a Galois extension $L/\Q$ with group $G$,
$L(s,S)=\zeta(s)L(s,\lambda,L/L^N)$ is the product of $\zeta$ and a nontrivial
quadratic Hecke $L$-function over $L^N$; it is holomorphic on
$\C\setminus\set1$, and $S$ is a nonnegative monomially certified weight in
the sense of \cref{grp:cor:detector}.
\end{proposition}

\begin{proof}
The induced character has degree $D$ and vanishes on unipotent and nonsplit
elements. All reflections of $N$ are involutions, which are split for
$q\equiv1\pmod4$, and their total $\lambda$-sum is zero. At a split
noninvolution the class meets $T_s$ in $t,t^{-1}$ and the centralizer has
order $a$; at the split involution the intersection with $T_s$ is a singleton
and the centralizer has order $2a$. In both cases the induction formula gives
$\delta(t)$.
\end{proof}

Thus \cref{grp:cor:detector}(i) applies to $S$: a one-sided bound for $R_S$
excludes zeros of the quadratic Hecke $L$-function. A multiplicity domination
of $\ord_\rho\zeta_L$ by $\ord_\rho L(s,S)$ is not established for odd $q$,
so the equivalence of \cref{grp:cor:detector}(ii) is not asserted.

\paragraph{Rich subgroups.}
A subgroup $H\le G$ is \emph{rich} if $\pi_H$ contains every irreducible
character, equivalently every irreducible representation has an $H$-fixed
vector; then $L^H$ is spectrally detecting. Let $h(q)$ be the maximal order of
a rich subgroup of $\PSL_2(q)$, so that the least degree of a single
spectrally detecting subfield is $\abs G/h(q)$. Rich subgroups, their
degree-sum bound and finite examples are studied in
\cite{BreuerHethelyiHorvathKulshammer2025Frobenius}.

\begin{theorem}[Single-subfield degrees over every odd field; \stU]
\label{grp:thm:odd-rich}
For every odd prime power $q\ge5$ other than $19$, $41$, $101$, $109$,
\[
 h(q)=
 \begin{cases}
 q-\sqrt q,&q=s^2,\ s\equiv3\pmod4,\\
 q-q^{3/4},&q\text{ a fourth power},\\
 (q+1)/2,&\text{otherwise},
 \end{cases}
 \qquad
 \min_{H\ \mathrm{rich}}[G:H]=\frac{q(q^2-1)}{2h(q)} ,
\]
while $h(19)=12$, $h(41)=24$ and $h(101)=h(109)=60$. In the first two cases of
the display an affine subgroup attains the maximum. For $q=3$, where
$G\simeq A_4$, $h(3)=2$. In particular $h(q)=(q+1)/2$ and the least degree is
$q(q-1)$, attained by the nonsplit torus, for every prime
$q\ge5$ outside $\set{19,41,101,109}$.
\end{theorem}

For example $h(121)=110$, $h(169)=85$ and $h(625)=500$, and $\PSL_2(109)$ has
a rich subgroup $A_5$ of order $60>55$.

\begin{proof}
\emph{Step 1: the nonsplit torus is rich.} Average each row of the table over
$T_n$. The trivial row, the Steinberg row and each exceptional row have
exactly one fixed dimension (on the nonsplit torus a cuspidal exceptional row
has the values $-\delta(t)$ away from the identity), and each ordinary
principal or cuspidal row has two, using $\sum_{t\ne1}\beta(t)=-1$. Hence
$h(q)\ge b$.

\emph{Step 2: $h(q)<q-1$.} If $H$ is rich, $\pi_H$ is nonnegative and contains
every irreducible, so by \cref{grp:thm:odd-gelfand} its degree is at least
$S(1)\ge D$ and $\abs H\le q-1$. Equality is impossible for
$q\equiv1\pmod4$, where the least degree is $D+1$. For $q\equiv3\pmod4$
equality would give $\pi_H=S$; at a generator $t$ of $T_n$, $S(t)=2$, so $H$
would contain a conjugate of $t$, but $b$ does not divide $q-1=2b-2$ for
$b>2$.

\emph{Step 3: reduction to an affine hyperplane.} Write $q=p^f$. By Dickson's
list \cite[Theorem~A.1]{GuralnickZieve2010Polynomials}, the exceptional groups
$A_4,S_4,A_5$ are excluded from the interval $(b,q-1)$ of orders when
$q\ge121$, since then $b>60$; smaller $q$ are treated in Step~6. A proper
subgroup of a torus normalizer has order at most $b$. A subfield group
$\PSL_2(t)$ with $q=t^j$ has order $(t^3-t)/2$, which exceeds $q-1$ for $j=2$
and is less than $b$ for $j\ge3$; a subfield group $\PGL_2(t)$ requires even
$j$, and its order $t^3-t$ exceeds $q-1$ for $j=2$ and is less than $b$ for
$j\ge4$. For the affine types, a nontrivial normal elementary abelian
$p$-subgroup has a unique common fixed point on the projective line, which
its normalizer preserves, so the group lies in a Borel subgroup; the
prime-to-$p$ subgroups of affine type are cyclic of order at most $(q-1)/2$.
It remains to consider $H\le B=U\rtimes T_s$ with $U=(\Fq,+)$ and
$T_s=(\Fq^\times)^2$ acting by multiplication. Let $A=H\cap U$ and let $C$ be
the image of $H$ in $T_s$. If $A=U$ then $\abs H\ge q$, and if $A=0$ then
$\abs H\le(q-1)/2$. If $C=1$, a proper additive subgroup has order at most
$q/p$. Otherwise lift a generator of $C$ to an affine map and translate its
fixed point to zero, so that $H=A\rtimes C$. Put $\F_r=\F_p(C)$; then $A$ is
$\F_r$-linear and $\abs A=q/r^m$ with $m\ge1$. The condition $\abs H>b$ gives
$\abs C/r^m>1/2$ while $\abs C\mid r-1$, so $m=1$ and $C=\F_r^\times$: $A$ is
an $\F_r$-hyperplane. Since $C$ consists of squares in $\Fq^\times$, the
degree $k=[\Fq:\F_r]$ is even; in particular $q\equiv1\pmod4$.

\emph{Step 4: the exceptional constituent test.} We claim that
$A\rtimes\F_r^\times$ is rich if and only if $4\mid(q-1)/(r-1)$. Identify the
characters of $U$ with $\Fq$ through the trace pairing. The annihilator
$A^\perp$ is an $\F_r$-line, and its nonzero elements form one free
$C$-orbit, contained in one of the two $T_s$-orbits of nontrivial additive
characters. Whenever a representation contains a weight in this free orbit,
it has an $H$-invariant: sum the $C$-translates of a nonzero vector of that
weight; the summands have distinct $U$-weights and are $A$-fixed. Every
ordinary principal, ordinary cuspidal and Steinberg row contains all
nontrivial $U$-weights, their restrictions to $U$ being $\Reg_U+\triv$,
$\Reg_U-\triv$ and $\Reg_U$ by the unipotent column. The exceptional pair
$E_+,E_-$ consists of the two constituents of $\Ind_B^G\delta$, with $\delta$
the quadratic character of $T_s$ inflated to $B$. Each has a $\delta$-line on
restriction to $B$, their combined restriction to $U$ is $\Reg_U+\triv$, and so
each has exactly one trivial $U$-line and its other weights fill one of the
two $T_s$-orbits, these having size $(q-1)/2$ and the degree being $(q+1)/2$.
The constituent meeting $A^\perp\setminus\set0$ has an invariant by the
free-orbit argument; in the other the $A$-invariants are exactly its
$\delta$-line. Hence $H$ is rich if and only if $\delta|_C=1$. The kernel of
$\delta$ in $\Fq^\times$ is the subgroup of fourth powers, which gives the
claim.

\emph{Step 5: maximization.} For admissible $r$ the affine subgroup has order
$q(1-1/r)$, so one takes the smallest admissible even $k$ with $q=r^k$. If
$r\equiv3\pmod4$, then $1+r+\dots+r^{k-1}$ is divisible by four for every even
$k$; if $r\equiv1\pmod4$, divisibility is equivalent to $4\mid k$. Thus $k=2$
works exactly in the first case of the display. If $4\mid f$ then $k=4$
works, and $\sqrt q$ is itself an odd square. If $f$ is odd there is no even
$k$. If $f\equiv2\pmod4$ and $k=2$ fails, then $p\equiv1\pmod4$ and every even
$k\mid f$ is twice an odd number, so none works. The admissible maximizing
affine orders exceed $b$ whenever the corresponding case occurs ($q\ge9$).

\emph{Step 6: the finite exceptional check.} Let $h_0(q)$ be the value of the
display. For odd prime powers $5\le q<121$ retain the groups
$K\in\set{A_4,S_4,A_5}$ on Dickson's list with $h_0(q)<\abs K<q-1$; here $A_4$
always occurs, $S_4$ requires $q\equiv\pm1\pmod8$ and $A_5$ requires
$q\equiv0,\pm1\pmod5$. This leaves the following cases, with the value of
$\min_{\chi\in\grpIrr(G)}\dim V_\chi^K$:
\[
\begin{array}{c|rrrrrrrrrr}
q&17&19&31&41&71&79&81&89&101&109\\\hline
K&A_4&A_4&S_4&S_4&A_5&A_5&A_5&A_5&A_5&A_5\\
\min&0&1&0&1&0&0&0&0&1&1
\end{array}
\]
The entries are computed as follows. The numbers of nonidentity elements of
orders $(2,3,4,5)$ are $(3,8,0,0)$, $(9,8,6,0)$, $(15,20,0,24)$ in
$A_4,S_4,A_5$. For an order $n$ prime to $p$ the torus is split if
$n\mid(q-1)/2$ and nonsplit otherwise; all cyclic subgroups of that order are
conjugate in $G$, and each contains its full set of generators in $K$. For a
torus parameter of exponent $i$, the average of its linear character over
these generators is $\mu(e)/\varphi(e)$ with $e=n/\gcd(i,n)$. Thus the average
on the class of elements of order $n$ in $K$ is twice this number for a
matching principal row, its negative double for a matching cuspidal row, half
the corresponding value for a matching exceptional row, and zero for a row of
the opposite series. Adding the degree to the sum of the order counts times
the row averages and dividing by $\abs K$ gives the fixed dimension. The only
entry in the defining characteristic is $q=81$, $K=A_5$, $n=3$: the twenty
such elements are conjugate in $K$ and lie in one unipotent class of $G$,
where the two exceptional values are $(1\pm9)/2$ and the other rows have the
values $1,-1,0$. The computation, in rational arithmetic, gives exactly the
four stated exceptions; its output is\\
{\small\artifact{grh_transfers_density_gaps_and_zero_tests/verification/check_small_subgroups.txt}}.\\
As a separate check, the two unipotent types are counted in nine candidate
affine subgroups; for $q=625$ the subgroup with $r=5$ fixes dimensions $1,2$
in the exceptional pair, whereas its larger competitor with $r=25$ fixes
dimensions $0,1$. The output is\\
{\small\artifact{grh_transfers_density_gaps_and_zero_tests/verification/check_affine.txt}}.

\emph{Step 7: $q=3$.} The sum of the three linear characters of $A_4$ and its
irreducible character of degree three takes the values $6,2,0$ at the
identity, the involutions and the elements of order three, and is the unique
minimizing nonnegative character by the positive-mean argument. A subgroup of
order two fixes every linear character and one dimension in the remaining
representation, where an involution has trace $-1$. The lower bound $6$ for
the degree excludes larger rich subgroups.
\end{proof}

\paragraph{Families of subfields and unconditional detection.}
Let $B=U\rtimes T_s$ be a Borel subgroup, with $U\simeq(\Fq,+)$ of order $q$, and let
$N_n=N_G(T_n)$, dihedral of order $q+1$. Write $F_H=L^H$ for the fixed field of $H\le G$ in a
Galois extension $L/k$ of number fields with group $G$, and denote the irreducible
characters as follows: $\chi_\alpha=\Ind_B^G\alpha$ of degree $q+1$ for a character $\alpha$ of
$T_s$ with $\alpha\ne\alpha^{-1}$, taken up to inversion; $\theta_\beta$ of degree $q-1$ for a
character $\beta$ of $T_n$ with $\beta\ne\beta^{-1}$, up to inversion; $\triv$; $\grpSt$; and the
two exceptional characters $E_\pm$. For $q\equiv1\pmod4$ let $z$ be the involution of $T_s$ and
$\delta$ its quadratic character; then $E_\pm$ have degree $(q+1)/2$ and
$E_++E_-=\Ind_B^G\delta$. Put
\[
 \Pi_0=\sum_{\alpha(z)=1}\chi_\alpha,\qquad
 I=\frac{1+\delta(z)}2=\begin{cases}1,&q\equiv1\pmod8,\\0,&q\equiv5\pmod8.\end{cases}
\]

\begin{lemma}[Fixed spaces; \stU]\label[lemma]{grp:lem:odd-fixed}
Let $q\equiv1\pmod4$, $q\ge5$. The multiplicities of the irreducible characters in $\pi_U$,
$\pi_{N_n}$ and $\pi_B$ are
\[
\begin{array}{l|rrrrr}
 &\triv&\grpSt&\chi_\alpha&\theta_\beta&E_\pm\\ \hline
\text{degree}&1&q&q+1&q-1&(q+1)/2\\
\pi_U&1&1&2&0&1\\
\pi_{N_n}&1&1&1+\alpha(z)&1&I\\
\pi_B&1&1&0&0&0
\end{array}
\]
In particular $\pi_U+\pi_{N_n}$ contains every irreducible character.
\end{lemma}
\begin{proof}
The values are those recalled before \cref{grp:thm:odd-gelfand}; the exceptional
characters take the value $\delta(t)$ at split $t$, vanish at nonsplit elements, and take the
values $(1\pm\sqrt q)/2$ on the two unipotent classes, interchanged for the other one. $U$ has
$(q-1)/2$ nontrivial elements of each unipotent class, which gives the $U$ row; for $E_\pm$ the
sum is $(q+1)/2+(q-1)/2$. Since $b=(q+1)/2$ is odd, the $b$ involutions of $N_n\setminus T_n$ are
split. A principal character vanishes on $T_n\setminus\set1$ and equals $2\alpha(z)$ on these
involutions, so its average is $(q+1+2b\alpha(z))/(2b)=1+\alpha(z)$. A cuspidal character sums to
$q-1+2=q+1$ over $T_n$, because $\sum_{t\ne1}\beta(t)=-1$, and vanishes on the involutions. The
exceptional average is $((q+1)/2+b\delta(z))/(2b)=I$, and the Steinberg average is
$(q-(b-1)+b)/(2b)=1$. Finally $\pi_B=\triv+\grpSt$, the action on the projective line being
doubly transitive.
\end{proof}

\begin{theorem}[Unconditional two-field detection; \stU]\label{grp:thm:odd-twofield}
Let $q\equiv1\pmod4$, $q\ge5$, and let $L/k$ be a Galois extension of number fields with group
$G=\PSL_2(\Fq)$. Then
\begin{equation}\label{grp:eq:odd-identity}
 (q-1)\pi_{N_n}+\frac{q+1}2\pi_U-\Reg_G
 =(q-1)\Bigl(\triv+\frac12\pi_B+2\Pi_0+I(E_++E_-)\Bigr),
\end{equation}
and every summand on the right is induced from a linear character, with nonnegative integral
total coefficients. Consequently, at every $s_0\ne1$,
\[
 0\le\ord_{s_0}\zeta_{F_U},\ \ord_{s_0}\zeta_{F_{N_n}}\le\ord_{s_0}\zeta_L
 \le\frac{q+1}2\ord_{s_0}\zeta_{F_U}+(q-1)\ord_{s_0}\zeta_{F_{N_n}} ,
\]
the zeros of $\zeta_L$ away from $1$ are those of $\zeta_{F_U}$ together with those of
$\zeta_{F_{N_n}}$, and $\mathrm{RH}(\zeta_L)$ is equivalent to $\mathrm{RH}(\zeta_{F_U})$ and
$\mathrm{RH}(\zeta_{F_{N_n}})$. The fields have degrees $[F_U:k]=(q^2-1)/2$ and
$[F_{N_n}:k]=q(q-1)/2$. Moreover
\[
 \frac{\zeta_{F_U}^{(q+1)/2}\zeta_{F_{N_n}}^{q-1}}{\zeta_L}
 =\zeta_k^{q-1}\zeta_{F_B}^{(q-1)/2}\prod_{\alpha(z)=1}L(s,\alpha,L/F_B)^{2(q-1)}\,
 L(s,\delta,L/F_B)^{I(q-1)},
\]
the product over the $\alpha$ of $\Pi_0$, is holomorphic on $\C\setminus\set1$ with a pole of exact
order $3(q-1)/2$ at $1$.
\end{theorem}
\begin{proof}
By \cref{grp:lem:odd-fixed}, the left side of \eqref{grp:eq:odd-identity} has multiplicity
$(q-1)(1+\alpha(z))$ at $\chi_\alpha$, that is $2(q-1)$ in $\Pi_0$ and $0$ otherwise; $0$ at every
$\theta_\beta$; $(q-1)I$ at $E_\pm$; $(q-1)/2$ at $\grpSt$; and $3(q-1)/2$ at $\triv$. These are the
multiplicities on the right, where $\pi_B=\triv+\grpSt$. The characters $\chi_\alpha$ and
$E_++E_-$ are induced from nontrivial linear characters of $B$. Let $\Theta$ be the Heilbronn
character at $s_0$. By \eqref{grp:eq:heil}, $\langle\Theta,\Ind_H^G\lambda\rangle\ge0$ for every
linear $\lambda$, $\Res_C\Theta$ is an actual character on cyclic $C$, so
$\abs{\Theta(g)}\le\Theta(1)$, and averaging over $H$ gives
$0\le\ord_{s_0}\zeta_{F_H}\le\ord_{s_0}\zeta_L$. Pairing \eqref{grp:eq:odd-identity} with $\Theta$
gives the upper bound. Artin formalism turns \eqref{grp:eq:odd-identity} into the displayed
factorization, all local factors included; the Hecke factors are entire and nonzero at $1$,
and $\zeta_k$, $\zeta_{F_B}$ have simple poles there.
\end{proof}

For $q=13$ the fields have degrees $84$ and $78$, against $156$ for the best single
spectrally detecting field (\cref{grp:thm:odd-rich}), and
$\ord_{s_0}\zeta_L\le7\ord_{s_0}\zeta_{F_U}+12\ord_{s_0}\zeta_{F_{N_n}}$. Here $\Pi_0$ is the single
character attached to a cubic character $\alpha$ of $T_s\simeq C_6$, and $I=0$, so
$\zeta_{F_U}^7\zeta_{F_{N_n}}^{12}/\zeta_L=\zeta_k^{12}\zeta_{F_B}^6L(s,\alpha,L/F_B)^{24}$.

\begin{theorem}[Optimal maximum degree of a family; \stU]\label{grp:thm:odd-family}
Let $q\equiv1\pmod4$. The least possible largest degree in a spectrally detecting family of
subfields of a $\PSL_2(\Fq)$-extension is
\[
 \mathcal D(q)=\begin{cases}36,&q=9,\\(q^2-1)/2,&q\ne9.\end{cases}
\]
For $q\ne9$ it is attained by $(F_U,F_{N_n})$, which detects zeros unconditionally by
\cref{grp:thm:odd-twofield}. For $q=9$ it is attained by $F_{N_n}$ and $F_A$, of degrees $36$ and
$30$, with $A\simeq\PSL_2(3)\simeq A_4$ a subfield subgroup, and
\begin{equation}\label{grp:eq:q9}
 8\pi_{N_n}+10\pi_A+18\,\triv-\Reg_G=9\pi_B+23\pi_{H_+}+3\pi_{H_-},
\end{equation}
with $H_\pm\simeq A_5$ of index six, so that
$\ord_{s_0}\zeta_L\le8\ord_{s_0}\zeta_{F_{N_n}}+10\ord_{s_0}\zeta_{F_A}+18\ord_{s_0}\zeta_k$ at every
$s_0\ne1$ and $(k,F_{N_n},F_A)$ detects zeros unconditionally with largest degree $36$.
\end{theorem}
\begin{proof}
A family whose smallest subgroup has order $h_0$ has largest degree $\abs G/h_0$, and
\cref{grp:lem:odd-fixed} gives the families above. For $q\ge13$ let $\alpha$ be a faithful
character of the cyclic group $T_s$ and $P=\Ind_B^G\alpha$, irreducible of degree $q+1$. We show
$P^H=0$ for every $H$ with $\abs H>q$, so that a family of subgroups all of order greater than
$q$ misses $P$; this uses Dickson's list \cite[Theorem~A.1]{GuralnickZieve2010Polynomials}.

\emph{Borel type.} Conjugate $H$ into $B$ and let $A'=H\cap U$ and $C$ its image in $T_s$. If
$C=1$, $\abs H\le q$. Otherwise $H=A'\rtimes C$ after translation, $A'$ is a vector space over the
field $\F_r$ generated by $C$, and $A'\ne U$ gives $\abs H\le(q/r)(r-1)<q$. So $A'=U$ and $C\ne1$.
The sum of $P$ over $U$ is $(q+1)+(q-1)=2q$, and the $q$ lifts of each $c\in C\setminus\set1$ are
split and contribute $q(\alpha(c)+\alpha(c)^{-1})$; faithfulness gives
$\sum_{c\ne1}\alpha(c)=-1$, so the total is $0$.

\emph{Torus normalizers.} Subgroups of the split normalizer have order at most $q-1$; the only
subgroup of the nonsplit normalizer of order greater than $q$ is $N_n$, and $\alpha(z)=-1$ since
$(q-1)/2$ is even, so $\langle P,\pi_{N_n}\rangle=0$ by \cref{grp:lem:odd-fixed}.

\emph{Subfield type.} $\PSL_2(t)$ with $q=t^j$, $j\ge3$, and $\PGL_2(t)$ with $j\ge4$ have order
less than $q$, so let $q=t^2$, $t\ge5$. Every semisimple element of these groups is split in
$G$; they have $t(t+1)/2$ split and $t(t-1)/2$ nonsplit tori, every nonidentity semisimple
element lies in exactly one, the sum of $P$ over the nonidentity elements of each torus is
$-2$, and the $t^2-1$ nonidentity unipotent elements have $P=1$. The total is
$(t^2+1)+(t^2-1)-2t^2=0$.

\emph{Exceptional type.} With $\abs H>q\ge13$ and Dickson's existence conditions only
$(q,H)=(17,S_4)$ and $(25,A_5),(29,A_5),(41,A_5),(49,A_5)$ remain. An element of order $d$
prime to $p$ is split exactly when $d\mid(q-1)/2$; the average of $P$ over the generators of a
split cyclic group of order $d$ is $2\mu(d)/\varphi(d)$, it is $0$ on nonsplit elements and $1$ on
unipotent ones. With the element counts $(9,8,6,0)$ and $(15,20,0,24)$ of orders $(2,3,4,5)$ in
$S_4$ and $A_5$, the sums of $P$ are $18-18$, $26-30-20+24$, $30-30$, $42-30-12$ and $50-30-20$.

\emph{$q=5$.} Here $G\simeq A_5$; a character of degree three has sum zero over $A_4$, $D_{10}$ and
$D_6$, the only proper subgroups of order greater than five ($3-3$, $3-5+2$, $3-3$), so
$\mathcal D(5)=12$, attained by $U$ and $N_n$ of orders $5$ and $6$.

\emph{$q=9$.} Here $G\simeq A_6$. With $\pi_{H_\pm}=\triv+E_\pm$, labelled so that $E_+$ occurs in
$\pi_A$, the multiplicities of $\triv,\grpSt,P,\theta_1,\theta_2,E_+,E_-$ (degrees
$1,9,10,8,8,5,5$) are $(1,1,0,1,1,1,1)$ in $\pi_{N_n}$, $(1,1,1,0,0,2,0)$ in $\pi_A$ and
$(1,1,0,0,0,0,0)$ in $\pi_B$; the $A$ row follows from its three split involutions and eight
unipotent elements of order three, on which the exceptional values are $2$ and $-1$. These give
\eqref{grp:eq:q9} and the coverage. For optimality take a cuspidal character $\theta$ of degree
eight: subgroups of a Borel of order greater than ten contain $U$, on which $\theta$ has no
invariants; no torus normalizer has order greater than ten; the sums of $\theta$ over $A_4$ and
$S_4$ are $8-8$ and over $A_5$ are $8-20+24/2$, the elements of order five being nonsplit with
average cuspidal value $1/2$. Hence $\mathcal D(9)=\abs G/10=36$. The inequality follows from
\eqref{grp:eq:q9} as in \cref{grp:thm:odd-twofield}.
\end{proof}

\begin{proposition}[Three fields for every odd $q$; \stU]\label[proposition]{grp:prop:odd-three}
For every odd prime power $q\ge3$ and $G=\PSL_2(\Fq)$,
\[
 \Reg_G=\frac{q-1}2\bigl(\pi_{T_n}+\pi_B-2\,\triv\bigr)+\pi_U,\qquad
 \zeta_L\zeta_k^{q-1}=\zeta_{F_{T_n}}^{(q-1)/2}\zeta_{F_B}^{(q-1)/2}\zeta_{F_U} .
\]
Hence $\ord_{s_0}\zeta_L\le\frac{q-1}2(\ord_{s_0}\zeta_{F_{T_n}}+\ord_{s_0}\zeta_{F_B})
+\ord_{s_0}\zeta_{F_U}$ at every $s_0\ne1$, the three fields, of degrees $q(q-1)$, $q+1$ and
$(q^2-1)/2$, detect all zeros of $\zeta_L$ away from $1$, and $\mathrm{RH}(\zeta_L)$ is equivalent
to their three Riemann hypotheses.
\end{proposition}
\begin{proof}
By Step~1 of the proof of \cref{grp:thm:odd-rich}, $\triv$, $\grpSt$ and each exceptional
character have one $T_n$-fixed dimension and every other irreducible has two. The $U$-fixed
dimensions are one for $\grpSt$, two for $\chi_\alpha$, zero for $\theta_\beta$, and one or zero for
the exceptional pair according as $q\equiv1$ or $3\pmod4$, when their degrees are $(q+1)/2$ or
$(q-1)/2$. Comparing with the degrees gives the identity; for $q=3$ the exceptional pair is the
pair of nontrivial linear characters of $A_4$. Artin formalism gives the zeta identity, and the
local positivity argument of \cref{grp:thm:odd-twofield} gives the rest.
\end{proof}

The identities \eqref{grp:eq:odd-identity} and of \cref{grp:prop:odd-three} are checked in class
values and multiplicities for all prime powers up to $1000$ in the relevant classes, together
with the exceptional subgroup cases and \eqref{grp:eq:q9}, in\
{\small\artifact{ProveIt_research_2026-11-27/02_psl2_detectors/verification_psl2.json}};\
the proofs above do not depend on these checks.

\begin{remark}[Relation to zeros]\label[remark]{grp:rem:odd-zeros}
Spectral coverage concerns which factors occur in the Artin factorization.
Under Artin holomorphy in the open strip, all local zero orders are
nonnegative and richness of $T_n$ gives the equivalence of the Riemann
hypotheses for $\zeta_{L^{T_n}}$ and $\zeta_L$. Without that input a
meromorphic constituent may have negative order; the finite-group optimality
supplies no sign control and no bound on the smoothed sum over primes. For the
single field $L^{T_n}$ this cancellation is not excluded. For the families of
\cref{grp:thm:odd-twofield,grp:thm:odd-family,grp:prop:odd-three} it is excluded
unconditionally, by positive identities, as \cref{grp:thm:psl-transfer} does in
the even case.
\end{remark}

\subsection{A quartic field with group \texorpdfstring{$S_4$}{S4}}\label{grp:sec:s4}
{\rightskip=0pt plus 1fil\relax
\sources{grh_deficit_transfer_2026-11-25; grh_relative_prime_races_2026-11-25}
}

Let $f(X)=X^4-X-1$ and let $E$ be its splitting field. The discriminant of $f$
is $-283$. Modulo $2$ the polynomial is irreducible, having no root and not
being divisible by $X^2+X+1$; modulo $7$ it factors as
$(X-3)(X^3+3X^2+2X+5)$ with the cubic having no root. So the Galois group
contains a $4$-cycle and a $3$-cycle; its order is divisible by $12$, and it
is not $A_4$, which contains no $4$-cycle. Hence $\Gal(E/\Q)=S_4$, and for
$p\ne283$ the factorization type of $f$ modulo $p$ is the cycle type of
$\Frob_p$. We use the fields
\[
 F_0=E^{A_4}=\Q(\sqrt{-283}),\qquad F_1=E^{V_4},\qquad F_4=E^{S_3},\qquad
 F_6=E^{C_4},
\]
where $F_0$ is the maximal abelian subfield, $F_1$ is the splitting field of
the cubic resolvent $Y^3+4Y-1$, $F_4$ is the quartic field fixed by a point
stabilizer and $F_6$ is fixed by the cyclic group generated by a $4$-cycle.
Dirichlet GRH (\cref{src:thm:dgrh}) gives the Riemann hypothesis for
$\zeta_{F_0}$. Let $\chi_3$ be the standard character,
$\chi_{3'}=\chi_3\,\grpsgn$, and $\chi_2$ the two-dimensional character
factoring through $S_4/V_4$. With $\pi_4,\pi_6$ the permutation characters of
$F_4,F_6$, put $\phi=\pi_4+\pi_6$, let $M(gA_4)$ be the maximum of $\phi$ on
the coset $gA_4$, and $u_{\mathrm{def}}=M-\phi$:
\begin{center}\small
\begin{tabular}{lrrrrr}
\toprule
cycle type & $1^4$ & $1^2\,2$ & $2^2$ & $1\,3$ & $4$\\\midrule
class size & 1&6&3&8&6\\
$\chi_2$ & 2&0&2&$-1$&0\\
$\chi_3$ & 3&1&$-1$&0&$-1$\\
$\chi_{3'}$ & 3&$-1$&$-1$&0&1\\
$\pi_4=\triv+\chi_3$ &4&2&0&1&0\\
$\pi_6=\triv+\chi_2+\chi_{3'}$ &6&0&2&0&2\\
$\phi=\pi_4+\pi_6$ &10&2&2&1&2\\
$M$ &10&2&10&10&2\\
$u_{\mathrm{def}}=M-\phi$ &0&0&8&9&0\\
\bottomrule
\end{tabular}\end{center}
The decompositions of $\pi_4$ and $\pi_6$ can be checked by scalar products
with the five irreducible characters. The function $u_{\mathrm{def}}$ is the
least nonnegative deficit of \cref{nab:sec:deficit} for the two fields
$F_4,F_6$ with $b_1=b_2=1$, and its mean is
$\delta=(3\cdot8+8\cdot9)/24=4$.

All three nonlinear Artin $L$-functions of $S_4$ are entire: $\chi_2$ is
induced from a nontrivial character of $A_4/V_4\simeq C_3$, and
$\chi_3,\chi_{3'}$ are induced from the two extensions, to its stabilizer
$V_4\rtimes C_2$ of order $8$, of a nontrivial character of $V_4$; the
induced representations are irreducible of dimension three, by restriction to
the three characters of $V_4$ permuted transitively, and differ by the sign
twist. Hence $\zeta_{F_4}$ and $\zeta_{F_6}$ together detect every nonlinear
factor without cancellation.

The two theorems below come from two different constructions: a nonnegative
deficit for the fields $F_4,F_6$, tested by a finite difference at the scales
$x$ and $2x$, and a signed compressed race along the tower
$F_0\subset F_1\subset E$, tested at the scales $x$ and $x^2$. Their weights
are compared in \cref{grp:rem:s4-compare}.

\begin{theorem}[Two-class deficit test; \stR]\label{grp:thm:s4}
Assume Dirichlet GRH (\cref{src:thm:dgrh}). For a prime $p\ne283$ let
$u_p=8$ if $f$ has two irreducible quadratic factors modulo $p$, $u_p=9$ if it
has a linear and an irreducible cubic factor, and $u_p=0$ otherwise, and let
$K_2$ be the kernel $K_\lambda$ of \cref{ana:sec:fd} with $\lambda=2$. The Riemann
hypothesis for $\zeta_E$ is equivalent to the one-sided inequality
\[
 D(x)=\frac1x\sum_{\substack{p\le2x\\p\ne283}}u_p\log p\,K_2(p/x)-4\log2
 \le C_\eps x^{-1/2+\eps}\qquad(\eps>0,\ x\text{ sufficiently large}),
\]
and to its averaged variant
$\int_X^{2X}\grppos{D(x)}\,\dd x=O_\eps(X^{1/2+\eps})$ for every $\eps>0$. The
unweighted version, with $u_p=1$ on the same two types and $u_p=0$ elsewhere,
also works, with the centring $\tfrac{11}{24}\log2$ in place of $4\log2$.
\end{theorem}

\begin{proof}
For the weighted version apply \cref{nab:thm:deficit} with $c=2$,
$[S_4,S_4]=A_4$, $B=F_0$, the two fields $K_1=F_4$, $K_2=F_6$ and $b_1=b_2=1$.
The table gives $w=u_{\mathrm{def}}$ and $\delta=4$, and the only ramified
prime is $283$. Given Dirichlet GRH, the inequality, or its averaged
variant, is therefore equivalent to the Riemann hypothesis for $\zeta_{F_4}$
and $\zeta_{F_6}$. Since $\zeta_{F_4}=\zeta L(s,\chi_3)$ and
$\zeta_{F_6}=\zeta L(s,\chi_2)L(s,\chi_{3'})$ and the three nonlinear factors
are entire, both fields detect all nonlinear factors, and the functional
equations and the product formula
$\zeta_E=\zeta L(s,\grpsgn)L(s,\chi_2)^2L(s,\chi_3)^3L(s,\chi_{3'})^3$, with
$\zeta L(s,\grpsgn)=\zeta_{F_0}$, give the Riemann hypothesis for $\zeta_E$.
Conversely the Riemann hypothesis for $\zeta_E$ implies it for both subfields,
because $\zeta_E/\zeta_{F_4}$ and $\zeta_E/\zeta_{F_6}$ are entire by the
Aramata--Brauer theorem \cite{Brauer1947Zeta,MurtyMurty1997Nonvanishing}.

The unweighted version is the case of the regular action,
$\phi=\grpone_{\set1}$, described after \cref{nab:thm:deficit} (equivalently
\cref{nab:thm:nonsplit}\textup{(v)} with $N=A_4$): the two types constitute
$A_4\setminus\set1$, of density $1/2-1/24=11/24$.
\end{proof}

Term by term, $D(x)=P_u(2x)-P_u(x)-\delta\log2$ in the notation of
\cref{nab:sec:interface} (\cref{ana:sec:fd}).

\begin{theorem}[Single quartic prime race; \stR]\label{grp:thm:s4-race}
Assume Dirichlet GRH (\cref{src:thm:dgrh}). For a prime $p\ne283$ let
$w_p=8$ if $f$ splits completely modulo $p$, $w_p=-1$ if it has a linear and
an irreducible cubic factor, and $w_p=0$ otherwise; that is,
$w=8\,\grpone_{1^4}-\grpone_{1\,3}$. Then $\zeta_E$ satisfies the Riemann
hypothesis if and only if for every $\eps\in(0,1/2)$ there is $C_\eps$ with
\[
 P_w(x)-P_w(x^2)\le C_\eps x^{-1/2+\eps}
\]
for all sufficiently large integers $x$, where $P_w$ is as in
\cref{nab:sec:interface} with $S=\set{283}$.
\end{theorem}

\begin{proof}
Apply \cref{nab:thm:tower} to the tower $F_0\subset F_1\subset E$ of fields
Galois over $\Q$, with $H_0=A_4$, $H_1=V_4$, $H_2=\set1$, $n_1=3$, $n_2=4$ and
coefficients $(c_1,c_2)=(1,2)$. On the classes $1^4,1^2\,2,2^2,1\,3,4$ the
weights of \cref{nab:lem:local} are $w_1=3\,\grpone_{V_4}-\grpone_{A_4}
=(2,0,2,-1,0)$ and $w_2=4\,\grpone_{\set1}-\grpone_{V_4}=(3,0,-1,0,0)$, so
$w_1+2w_2=(8,0,0,-1,0)=w$, in accordance with \eqref{nab:eq:compressed}. Up
to a function holomorphic on $\Re s>1/2$,
\[
 D^{\mathrm{pr}}_w(s)=-\frac12\frac{Q_1'}{Q_1}(s)-\frac13\frac{Q_2'}{Q_2}(s),\qquad
 Q_1=\frac{\zeta_{F_1}}{\zeta_{F_0}},\quad Q_2=\frac{\zeta_E}{\zeta_{F_1}} .
\]
The support of $w$ has density $9/24=3/8$. By \cref{nab:thm:tower} the
inequality is equivalent to the relative Riemann hypothesis for
$\zeta_E/\zeta_{F_0}$, and since $F_0$ is quadratic, Dirichlet GRH gives the
Riemann hypothesis for $\zeta_{F_0}$, hence the equivalence with the Riemann
hypothesis for $\zeta_E$.
\end{proof}

\begin{remark}[Comparison of the weights]\label[remark]{grp:rem:s4-compare}
The deficit of \cref{grp:thm:s4} is tested with the orientation opposite to
that of the race: an upper bound for $P_u(2x)-P_u(x)-\delta\log2$ is an upper
bound for $R_{-u}(x)-R_{-u}(2x)$, and the series of $-u_{\mathrm{def}}$ has the
positive coefficients $b_j$ at $-\zeta_{F_j}'/\zeta_{F_j}$ (proof of
\cref{nab:thm:deficit}). The weights to be compared are therefore
$-u_{\mathrm{def}}$ and $w_{\mathrm{race}}=w$ of \cref{grp:thm:s4-race}. Their
values, and those of their difference, are as follows.
\begin{center}\small
\begin{tabular}{lrrrrr}
\toprule
cycle type & $1^4$ & $1^2\,2$ & $2^2$ & $1\,3$ & $4$\\\midrule
$u_{\mathrm{def}}$ (\cref{grp:thm:s4}) &0&0&8&9&0\\
$w_{\mathrm{race}}$ (\cref{grp:thm:s4-race}) &8&0&0&$-1$&0\\
$w_{\mathrm{race}}-(-u_{\mathrm{def}})$ &8&0&8&8&0\\
$8\cdot\grpone_{A_4}=4(\triv+\grpsgn)$ &8&0&8&8&0\\
\bottomrule
\end{tabular}\end{center}
So $w_{\mathrm{race}}=-u_{\mathrm{def}}+4(\triv+\grpsgn)$ on every class. The
difference $4(\triv+\grpsgn)$ is a function of the image of $\Frob_p$ in
$\Gal(F_0/\Q)$, that is of $\chi_{-283}(p)$, and its $L$-functions are $\zeta$
and $L(s,\chi_{-283})$, whose product is $\zeta_{F_0}$. This is the term the
Dirichlet-GRH input removes in each theorem: in \cref{grp:thm:s4} it removes
the abelian part $M=6\,\triv+4\,\grpsgn$ of the deficit, and in
\cref{grp:thm:s4-race} it removes $\zeta_{F_0}$ at the bottom of the tower.

The table is an identity of class functions only. The two theorems test the
two weights with different statistics, an increment at the scales $x,2x$ with
a main term $4\log2$, and a constant-free difference at the integer scales
$x,x^2$; the identity does not convert one inequality into the other. That
both are equivalent to the Riemann hypothesis for $\zeta_E$, given Dirichlet
GRH, is the content of the two theorems. In what an arithmetic construction
would have to count they differ: the deficit has nonnegative weights, support
of density $11/24$, and never observes completely split primes; the race has
mean zero, one positive and one negative class, support of density $9/24$,
and needs the completely split primes with the large weight $8$. The
unweighted form in \cref{grp:thm:s4} tests $\zeta_E$ through the regular
action, with the single coefficient $1/24$. The Sylow-filtered race of
\cref{nab:thm:sylow} for $S_4$, with weights $17,-3,-1$ on $1^4,2^2,1\,3$
(\cref{nab:ex:sylowsn}), is supported on all of $A_4$.

None of these weights is a function of $p$ modulo a fixed modulus
(\cref{nab:thm:periodic}): congruences determine the sign of $\Frob_p$ through
$\chi_{-283}$ but cannot separate the identity from the $3$-cycles, both of
which lie in $A_4$.
\end{remark}

\begin{remark}[Minimal support of strict-sign deficits]\label[remark]{grp:prop:s4-support}
Let $f\ge0$ be a class function on $S_4$ whose inner products with the three nonlinear
irreducible characters, the new characters for $N=A_4$, are strictly negative. Then
\eqref{nab:eq:identity} gives $f(u)>f(1)\ge0$ on $A_4\setminus\set1$, as in
\cref{nab:thm:minimax}(b), whose normalization $f(1)=0$ is not needed here; so the union
of the classes $2^2$ and $1\,3$ is the unique minimal support. For a function with
weights $u$ on $2^2$ and $v$ on $1\,3$ and zero elsewhere, the exact condition is $u>0$
and $v>3u/4$: the inner products with $\chi_3$ and $\chi_{3'}$ are both $-3u/24$, and
that with $\chi_2$ is $(6u-8v)/24$. The minimality refers to this strict-sign
architecture and is not a lower bound for all possible criteria. Both deficit weights of
\cref{grp:thm:s4} satisfy the condition, with $(u,v)=(8,9)$ and $(1,1)$.
\end{remark}

\subsection{One splitting test for six quintic fields}\label{grp:sec:six}
{\rightskip=0pt plus 1fil\relax
\sources{grh-exponent-six-and-transfers}
}

Fix $D=-12451$ or $D=-37363$, and put $E=\Q(\sqrt D)$ and $B=E(\zeta_5)$. For
these two discriminants $\Cl(E)[5]\simeq\F_5^2$, and the six lines of this
plane give six pairwise nonisomorphic totally real quintic fields
$K_1,\dots,K_6$, the reflected fields of \cref{dir:sec:reflected}
(\cref{dir:prop:reflection}); they are the six fields of \cref{dir:prop:ranktwo}, with
the defining polynomials of \cref{dir:tab:ranktwo}.

\paragraph{The joint Galois group.}
Choose virtual units $\alpha_1,\alpha_2$ representing a basis of
$\Cl(E)[5]$. Over $B$ their radical classes are independent, by restriction
followed by norm, so adjoining both fifth roots gives an extension $M/B$ with
group $V=\F_5^2$ by Kummer theory. Complex conjugation $c$ is central in
$\Gal(M/\Q)$. The six line evaluations of $V$ separate points, and each
quintic splitting field detects the full scalar group; hence the joint normal
closure $T$ of $K_1,\dots,K_6$ is $M^{\langle c\rangle}$, and
\[
 \Gal(T/\Q)=V\rtimes\F_5^\times,\qquad H:=T^V=B^{\langle c\rangle},
\]
the scalars acting on both coordinates. On the five embeddings of the field
attached to a line, an element $(v,a)$ acts by $j\mapsto aj+\lambda(v)$, with
$\lambda$ a linear form representing the corresponding line of the dual
plane. At $p\nmid5D$ the scalar of Frobenius is
\[
 \mu(p)=\chi_D(p)\,(p\bmod5)\in\F_5^\times ,
\]
since Frobenius raises $\zeta_5$ to the power $p$ and acts on $E$ through
$\chi_D$, while simultaneous complex conjugation is the identity on $H$. For
$a\ne1$ the affine map is conjugate to multiplication by $a$; for $a=1$ a
nonzero translation is a $5$-cycle, and exactly one of the six linear forms
vanishes on a given $v\ne0$. Chebotarev's theorem gives the densities
\[
\begin{array}{lll}
 (v,a) & \text{residue degrees in the six fields} & \text{density}\\
 \hline
 v=0,\ a=1 & (1,1,1,1,1)\text{ in all six} & 1/100\\
 v\ne0,\ a=1 & (1,1,1,1,1)\text{ in one, }(5)\text{ in five} & 24/100\\
 a=-1 & (1,2,2)\text{ in all six} & 25/100\\
 a=2\text{ or }3 & (1,4)\text{ in all six} & 50/100
\end{array}
\]
unconditionally; each of the six choices of the uniquely splitting quintic in
the second row has density $4/100$. In particular the set of primes splitting
completely in two distinct fields $K_i,K_j$ does not depend on the pair.

\begin{theorem}[Six-zeta test; \stR]\label{grp:thm:six}
Assume Dirichlet GRH (\cref{src:thm:dgrh}). Let $S$ be a finite set of primes
containing those dividing $5D$, fix $i\ne j$, and put
\[
 W(p)=25\,\grpone_{\set{p\text{ splits completely in }K_i\text{ and }K_j}}
 -\grpone_{\set{\mu(p)=1}} ,
\]
which takes the values $24,-1,0,0$ in the four rows above and has mean zero.
The following are equivalent.
\begin{enumerate}[label=\textup{(\arabic*)}]
\item All six $\zeta_{K_i}$ satisfy the Riemann hypothesis.
\item $\zeta_T$ satisfies the Riemann hypothesis.
\item For every $\eps>0$ there is $C_\eps$ with
$P_W(x)-c_W\le C_\eps x^{-1/2+\eps}$ for all sufficiently large $x$, where
$c_W=\lim_{x\to\infty}P_W(x)$.
\end{enumerate}
Under these conditions $\abs{P_W(x)-c_W}\ll_{D,S}x^{-1/2}\log^2(2x)$.
\end{theorem}

\begin{proof}
Let $\tau_i$ be the permutation character on the five embeddings of $K_i$
minus the trivial character. The table gives $W=\sum_i\tau_i$ on all Frobenius
classes. The regular character of $\Gal(T/\Q)$ equals the inflated regular
character of $\Gal(H/\Q)$ plus $4\sum_i\tau_i$, the sum having the value $100$
at the identity and zero elsewhere. Hence
\[
 \zeta_T(s)=\zeta_H(s)F(s)^4,\qquad
 F(s)=\prod_{i=1}^6\frac{\zeta_{K_i}(s)}{\zeta(s)} .
\]
Since $H/\Q$ is abelian, $\zeta_H$ is a product of Dirichlet $L$-functions
(\cref{dir:prop:abelian}).
Dirichlet GRH makes each ratio in $F$ holomorphic on $\Re s>1/2$, with a
nonzero removable value at $1$, and makes $\zeta_H$ zero-free there. Zeros of
the six numerators in the right half of the strip therefore contribute
positive orders to $F$, and the functional equations give the equivalence of
(1) and (2).

Euler factors give $D^{\mathrm{pr}}_W=-F'/F-J$ with $J$ holomorphic on $\Re s>1/2$: at good
primes the terms with $p^j$, $j\ge2$, are bounded by $24(\log p)p^{-j\Re s}$,
and the corrections at the excluded primes are holomorphic on $\Re s>0$. This
is \eqref{grp:eq:decomp} with $Q_1=F$, so (3) implies that $F$ has no zero in
$\Re s>1/2$ by \cref{grp:prop:interface}(b), which gives (1). Conversely (2)
gives the two-sided bound by \cref{grp:prop:interface}(d).
\end{proof}

\begin{corollary}[Information missing from residue classes; \stU]
\label[corollary]{grp:cor:six-variance}
For every $m\ge1$ and every $h:(\Z/m\Z)^\times\to\C$, in the notation of
\cref{nab:thm:periodic} for $\Gal(T/\Q)$ with $q=m$,
\[
 \mathcal E_m(W,h)=6+\frac1{\varphi(m)}\sum_{a\in(\Z/m\Z)^\times}\abs{h(a)}^2 .
\]
\end{corollary}

This is \cref{nab:thm:periodic}: the commutator subgroup of $V\rtimes\F_5^\times$ is $V$
(commutators with the scalar $2$ give every translation), and $W$ has mean zero on each
coset of $V$, so $m_q=0$ and the minimum is $\norm W^2=(24^2+24)/100=6$.

Among the primes up to $5000$ outside $S$ and the polynomial discriminants,
the first prime splitting completely in all six fields is $269$ for
$D=-12451$ and $769$ for $D=-37363$
(\artifact{grh-exponent-six-and-transfers/data/splitting.json}).

\subsection{Kummer fields and the primitive-root target}\label{grp:sec:kummer}
{\rightskip=0pt plus 1fil\relax
\sources{grh-transfer; grh_deficit_transfer_2026-11-25; grh_kummer_window_transfer_2026-11-25}
}

Hooley's conditional proof of Artin's primitive-root conjecture for a base $a$
uses the Riemann hypothesis for the Dedekind zeta functions of the fields
$\Q(\zeta_m,a^{1/m})$ with $m$ squarefree \cite{Hooley1967Artin,Moree2012Artin}.
Already $m=3$ is nonabelian, so these fields are the first targets beyond
\cref{src:thm:dgrh} in this direction. This subsection gives the exact
detectors for them, the form of the tests in finite sums, a criterion for
finite blocks of prime indices, and the consequences for the base $2$.

\subsubsection{The fields and their prime weights}
Fix a nonzero integer $a$, not $-1$ and not a square. For squarefree $m\ge2$
let $B_m=\Q(\zeta_m)$, $E_m=B_m(a^{1/m})$ and $h_m=[E_m:B_m]$, a divisor of
$m$; $E_m/\Q$ is Galois, being the splitting field of $X^m-a$. For
$p\nmid am$ define the nonresidue weight and its density
\begin{equation}\label{grp:eq:kummer-weight}
 w_{a,m}(p)=\begin{cases}
 1,&p\equiv1\pmod m\text{ and }a^{(p-1)/m}\not\equiv1\pmod p,\\
 0,&\text{otherwise},
 \end{cases}
 \qquad
 \delta_{a,m}=\frac{1-h_m^{-1}}{\varphi(m)} .
\end{equation}
Splitting in $B_m$ is the congruence $p\equiv1\pmod m$; then $\F_p^\times$
contains the $m$th roots of unity and $X^m-a$ splits if and only if $a$ is an
$m$th power, which is tested by the modular exponent. So $w_{a,m}$ is the
indicator of the primes which split completely in $B_m$ but not in $E_m$. For
an odd prime $\ell$ with $a$ not an $\ell$th power in $\Q$, $h_\ell=\ell$ and
$\delta_{a,\ell}=1/\ell$: an $\ell$th root of $a$ in $B_\ell$ would, on taking
norms, make $a^{\ell-1}$ a rational $\ell$th power, hence $a$ an $\ell$th
power.

For the base $a=2$ and an odd prime $\ell$ put
\[
 F_\ell=\Q(\zeta_\ell),\qquad B_\ell=\Q(2^{1/\ell}),\qquad
 K_\ell=F_\ell B_\ell=\Q(\zeta_\ell,2^{1/\ell}).
\]
The binomial $X^\ell-2$ is Eisenstein at $2$, so $[B_\ell:\Q]=\ell$, and
$[K_\ell:\Q]=\ell(\ell-1)$ since $\ell$ and $\ell-1$ are coprime. The Galois
group acts on the roots as the full affine group $t\mapsto at+b$ of
$\F_\ell$, and for $p\nmid2\ell$ the dilation of $\Frob_p$ is $a=p\bmod\ell$.
If $a\ne1$ the affine map has a fixed point; if $a=1$ the identity splits the
polynomial completely and every nonidentity translation is an $\ell$-cycle.
Let $V_\ell$ be the permutation character minus the trivial character, with
values $\ell-1$ at the identity, $-1$ at nonidentity translations and $0$ for
$a\ne1$. It is irreducible, and $V_\ell=\Ind_N^G\psi$ for any nontrivial
character $\psi$ of the translation subgroup $N$. Artin induction gives
\begin{equation}\label{grp:eq:kummer-fact}
 L_\ell(s):=L(s,V_\ell)=L(s,\psi,K_\ell/F_\ell)
 =\frac{\zeta_{B_\ell}(s)}{\zeta(s)},\qquad
 \zeta_{K_\ell}(s)=\zeta_{F_\ell}(s)L_\ell(s)^{\ell-1},
\end{equation}
because the regular character of $G$ is the inflated regular character of
$G/N$ plus $(\ell-1)V_\ell$. The Hecke $L$-function $L_\ell$ is entire and
nonzero on $\Re s\ge1$. For $p\nmid2\ell$ let
\[
 I_\ell(p)=\grpone_{p\equiv1\ (\ell)},\qquad
 g_\ell(p)=\grpone_{p\equiv1\ (\ell),\ 2^{(p-1)/\ell}\equiv1\ (p)},\qquad
 f_\ell=I_\ell-g_\ell=w_{2,\ell},
\]
all set to zero at $p\in\set{2,\ell}$. The three affine possibilities give
\begin{equation}\label{grp:eq:kummer-ind}
 g_\ell=\frac{I_\ell+V_\ell}{\ell},\qquad
 f_\ell=\frac{(\ell-1)I_\ell-V_\ell}{\ell},
\end{equation}
where $V_\ell$ stands for its value at Frobenius; equivalently
$r_\ell(p)-1=V_\ell(\Frob_p)$ for the number $r_\ell(p)$ of roots of
$X^\ell-2$ modulo $p$. Logarithmic differentiation of the Euler factors, as in
the proof of \cref{grp:cor:detector}, gives for $\Re s>1$
\begin{align}
 D^{\mathrm{pr}}_{g_\ell}(s)&=-\frac1{\ell(\ell-1)}\frac{\zeta_{F_\ell}'(s)}{\zeta_{F_\ell}(s)}
 -\frac1\ell\frac{L_\ell'(s)}{L_\ell(s)}+A_\ell(s),\label{grp:eq:kummer-pg}\\
 D^{\mathrm{pr}}_{f_\ell}(s)&=-\frac1\ell\frac{\zeta_{F_\ell}'(s)}{\zeta_{F_\ell}(s)}
 +\frac1\ell\frac{L_\ell'(s)}{L_\ell(s)}+A'_\ell(s),\label{grp:eq:kummer-pf}
\end{align}
with $A_\ell,A'_\ell$ holomorphic on $\Re s>1/2$; the residues at $1$ are
$\delta_{g_\ell}=1/(\ell(\ell-1))$, $\delta_{f_\ell}=1/\ell$ and
$\delta_{I_\ell}=1/(\ell-1)$.

The orientation is essential. By \eqref{grp:eq:kummer-pf} the nonresidue sum
$R_{f_\ell}$ has a \emph{negative} residue at a real zero of $L_\ell$: an
upper bound for the surplus of nonresidue primes is not a detector. What
detects zeros is an upper bound for the split-prime sum $R_{g_\ell}$ or,
given Dirichlet GRH, an upper bound for the nonresidue \emph{deficit}
$-R_{f_\ell}$. The source inequality is an upper bound for a surplus of
nonsplit quadratic primes; substituting a nonresidue weight for its quadratic
weight gives the wrong sign.

\subsubsection{Completely split primes and tests in finite sums}

For $M/\Q$ Galois of degree $n$ let $R_M=R_w$ for the weight
$w=\grpone_{\set1}$ on $\Gal(M/\Q)$, so that $\delta_w=1/n$ and $R_M$ is the
centred ramp sum over the unramified primes that split completely in $M$. It is
exactly $1/n$ times the remainder of \cref{nab:prop:dedekind} for $K=E=M$, so by
\cref{grp:prop:split} and \cref{nab:prop:dedekind}(b) the Riemann
hypothesis for $\zeta_M$ is equivalent to $\sup_{x\ge2}\sqrt x\,R_M(x)<\infty$,
and then $\abs{R_M(x)}=O_M(x^{-1/2})$. The exact value of
$\limsup\sqrt x\,R_M(x)$ under the Riemann hypothesis for $\zeta_M$ is
\cref{zer:thm:endpoint}.

Let now $E\subset M$ with $E/\Q$ abelian, let $I$ be the set of primes,
unramified in $M$, which split completely in $E$ but not in $M$, and
$\delta=1/[E:\Q]-1/[M:\Q]$. After centring, $R_I=R_E-R_M$ exactly, and by
\cref{nab:thm:nonsplit}\textup{(i)}$\Leftrightarrow$\textup{(ii)}, applied to
$\Gal(M/\Q)$ and its normal subgroup $\Gal(M/E)$, the \emph{lower} bound
$R_I(x)\ge-Cx^{-1/2}$ is equivalent to the Riemann hypothesis for $\zeta_M$. Its
centring constant is removed by the dyadic increment
\begin{equation}\label{grp:eq:BI}
 B_I(x)=R_I(2x)-R_I(x)
 =\frac1{2x}\sum_{p\in I,\ p\le x}\log p
 +\sum_{\substack{p\in I\\x<p\le2x}}\log p\Bigl(\frac1p-\frac1{2x}\Bigr)
 -\delta\log2 .
\end{equation}
Under the Riemann hypothesis for $\zeta_M$, $\abs{B_I(x)}=O_M(x^{-1/2})$.
Conversely, if $B_I(x)\le Cx^{-1/2}$, telescoping on $x,2x,4x,\dots$ gives
$-R_I(x)\le Cx^{-1/2}/(1-2^{-1/2})$, since the remainder tends to zero; and it
is enough to check integers $x\ge2$, because $P_I$ is continuous with
derivative $\theta_I(x)/x^2=O(1/x)$ off the primes. Thus $\zeta_M$ satisfies the
Riemann hypothesis if and only if $\sup_{m\in\Z,\ m\ge2}\sqrt m\,B_I(m)<\infty$: a form
of \cref{nab:thm:nonsplit}\textup{(v)} for the kernel with $c=2$, with the exponent
$-1/2$ exactly.

\begin{example}[The splitting field of $X^\ell-2$; \stR]\label[example]{grp:thm:kummer-single}
Assume Dirichlet GRH (\cref{src:thm:dgrh}). Let $\ell$ be an odd prime and
$K_\ell=\Q(\zeta_\ell,2^{1/\ell})$. Let $I$ be the set of primes $p\nmid2\ell$
for which $X^\ell-2$ is irreducible modulo $p$; equivalently
$p\equiv1\pmod\ell$ and $2^{(p-1)/\ell}\not\equiv1\pmod p$. By the description of the
affine action, $I$ is exactly the set of good primes splitting completely in $F_\ell$ but
not in $K_\ell$, of density $1/(\ell-1)-1/(\ell(\ell-1))=1/\ell$; here $\Gal(K_\ell/F_\ell)$
is the translation subgroup, the commutator subgroup of the affine group. The
preceding paragraph with $E=F_\ell$, $M=K_\ell$ and $\delta=1/\ell$ gives
\[
 \text{the Riemann hypothesis for }\zeta_{K_\ell}
 \quad\Longleftrightarrow\quad
 \sup_{m\in\Z,\ m\ge2}\sqrt m\,B_I(m)<\infty .
\]
\end{example}

For $\ell=3$ the group is $S_3$. The three Frobenius types are completely
split (density $1/6$), linear times irreducible quadratic ($1/2$) and
irreducible cubic ($1/3$). The middle condition is $p\equiv2\pmod3$, whose
centred sum already has the square-root estimate; the missing estimate
distinguishes the two types within $p\equiv1\pmod3$. For example $7$ and $31$
are both $7$ modulo $12$, but $2$ is a cube only modulo $31$. No fixed modulus
makes this distinction (\cref{nab:thm:periodic}): the maximal abelian subfield of
$K_\ell$ is $F_\ell$, so within every progression $a\bmod q$ with $a\equiv1\pmod\ell$
the two splitting types have relative densities $1/\ell$ and $1-1/\ell$. Numerical values of $\sqrt x\,B_I(x)$ for
$\ell=3,5$ and $x\le10^5$ are in
\artifact{grh-transfer/checks/kummer-output.txt}.

\begin{corollary}[General base and squarefree index; \stR]\label[corollary]{grp:cor:kummer-general}
Assume Dirichlet GRH. Let $a$ and squarefree $m\ge2$ be as above. If for every
$\eps>0$
\[
 P_{w_{a,m}}(2x)-P_{w_{a,m}}(x)-\delta_{a,m}\log2\le C_{a,m,\eps}x^{-1/2+\eps}
\]
for all sufficiently large $x$, the sums omitting the primes dividing $am$,
then $\zeta_{E_m}$ satisfies the Riemann hypothesis; the same holds under the
corresponding mean bound for the positive part. The converse holds as well.
\end{corollary}

\begin{proof}
This is the discussion before \cref{grp:thm:kummer-single} with $E=B_m$, $M=E_m$ and
$\delta=\delta_{a,m}$; in the notation of \cref{grp:prop:interface},
$D^{\mathrm{pr}}_{-w_{a,m}}=-\frac1{[E_m:\Q]}\zeta_{E_m}'/\zeta_{E_m}
+\frac1{\varphi(m)}\zeta_{B_m}'/\zeta_{B_m}$ up to a holomorphic function, and
the second term is holomorphic on $\set{\Re s>1/2}\setminus\set1$ by Dirichlet
GRH.
\end{proof}

This is a direct analogue of the nonsplit quadratic prime statistic of the
source proof: the congruence $p\equiv1\pmod m$ is retained, and the quadratic
character is replaced by the $m$th power-residue test. For $a=2$, $m=3$ the
test consists exactly of the primes $p\equiv1\pmod3$ for which $2$ is not a
cube.

\subsubsection{One averaged inequality for a block of prime indices}
For a real weight $w$ and $h>0$ write $R_w(v)$ for $R_w(e^v)$ in this
paragraph and put $\overline R_w(v)=h^{-1}\int_0^hR_w(v+u)\,\dd u$.

\begin{lemma}[Moving averages; \stU]\label[lemma]{grp:lem:moving}
Let $h>0$ and
\[
 J_h(s)=\frac{e^{h(s-1)}-1}{h(s-1)},\qquad J_h(1)=1 .
\]
For $\Re s>1$ the transform
$\int_0^\infty\overline R_w(v)e^{-(s-1)v}\,\dd v$ equals $J_h(s)$ times the
transform \eqref{grp:eq:mellin} of $R_w$, plus an entire function. The
multiplier $J_h$ has no zeros in $\Re s<1$ and is positive for real $s$.
\end{lemma}

\begin{proof}
Replace $v$ by $v+u$ in the inner integral; the omitted initial interval
$[0,u]$ contributes an entire function after integration over $u\in[0,h]$.
The multiplier is the average of $e^{u(s-1)}$ over $[0,h]$; its only zeros are
$s=1+2\pi ik/h$, $k\ne0$, and it is positive on the real axis.
\end{proof}

\begin{theorem}[Positive-block criterion; \stR]\label{grp:thm:block}
Fix a nonempty finite set $S$ of odd primes, positive real numbers $b_\ell$
($\ell\in S$) and $h>0$. For $1/2<\beta<1$, the inequality
\begin{equation}\label{grp:eq:block}
 \sum_{\ell\in S}b_\ell\,\ell\,\overline R_{g_\ell}(v)\le C_\beta e^{(\beta-1)v}
 \qquad(v\text{ sufficiently large})
\end{equation}
implies $\zeta_{K_\ell}(s)\ne0$ for $\Re s>\beta$, $s\ne1$, for all
$\ell\in S$. Under Dirichlet GRH the same follows from the nonresidue deficit
bound
\begin{equation}\label{grp:eq:nonres}
 -\sum_{\ell\in S}b_\ell\,\ell\,\overline R_{f_\ell}(v)\le C_\beta e^{(\beta-1)v}.
\end{equation}
If either inequality holds for every $1/2<\beta<1$, with constants depending
on $S,h,\beta$, all the fields $K_\ell$, $\ell\in S$, satisfy the Riemann
hypothesis for their Dedekind zeta functions. Conversely these Riemann
hypotheses imply \eqref{grp:eq:block} and \eqref{grp:eq:nonres} for every
$\beta>1/2$.
\end{theorem}

\begin{proof}
By \eqref{grp:eq:kummer-pg}, at a point $\rho$ with $\beta<\Re\rho<1$ the sum
of the ramp transforms in \eqref{grp:eq:block}, before averaging, has the
residue
\[
 \frac{\sum_{\ell\in S}b_\ell\bigl(\ord_\rho L_\ell
 +\frac1{\ell-1}\ord_\rho\zeta_{F_\ell}\bigr)}{\rho(1-\rho)} .
\]
All orders are nonnegative integers, so there is a pole whenever any factor
vanishes and no cancellation is possible, even for coincident nonreal zeros.
Averaging multiplies the residue by $J_h(\rho)\ne0$, which is positive at real
$\rho$ (\cref{grp:lem:moving}). The moving averages are continuous and
$O((1+v)^2)$. Hence \cref{ana:thm:onesided} excludes all such poles, and
\eqref{grp:eq:kummer-fact} gives the conclusion. For \eqref{grp:eq:nonres} use
\eqref{grp:eq:kummer-pf}: Dirichlet GRH removes the terms with
$\zeta_{F_\ell}$ in this region, and the remaining orders of $L_\ell$ enter
with the same sign as before. All factors are nonzero on $\Re s\ge1$, the
pole of the zeta functions excepted. Let $\beta\downarrow1/2$ and use the
functional equations.

Conversely, under the Riemann hypothesis for $\zeta_{K_\ell}$,
\cref{grp:prop:interface}(d) gives
$R_w(v)=O_\ell(e^{-v/2}(1+v)^2)$ for $w=g_\ell$ and $w=I_\ell$, which are in
the situation of \cref{grp:prop:interface}(b) by \eqref{grp:eq:kummer-pg} and
$D^{\mathrm{pr}}_{I_\ell}=-\frac1{\ell-1}\zeta_{F_\ell}'/\zeta_{F_\ell}$ up to a holomorphic
function, hence also for $f_\ell=I_\ell-g_\ell$; averaging over a fixed window
and summing finitely many terms preserves the bound.
\end{proof}

The pointwise bounds, without averaging, are stronger hypotheses and suffice
as well. Taking $h=\log2$ one bounds the mean of the centred ramp statistic
over an octave. A general smoothing kernel could introduce zeros that hide a
pole; $J_h$ has none in the open critical strip.

\begin{corollary}[Periodic obstruction for a block; \stU]\label[corollary]{grp:cor:periodic}
For a block $S$ as above let $n=\prod_{\ell\in S}\ell$ and
$Z(p)=\sum_{\ell\in S}b_\ell V_\ell(\Frob_p)$ at the good primes. For every
fixed modulus $q$ and every $a:(\Z/q\Z)^\times\to\C$, in the notation of
\cref{nab:thm:periodic} for $\Gal(K_n/\Q)$,
\[
 \mathcal E_q(Z,a)=\sum_{\ell\in S}b_\ell^2
 +\frac1{\varphi(q)}\sum_{b\in(\Z/q\Z)^\times}\abs{a(b)}^2 .
\]
When $q$ is divisible by all $\ell\in S$, the same residual mean square
$\sum_\ell b_\ell^2$ is the minimum for $\sum_\ell b_\ell\,\ell\,g_\ell$.
\end{corollary}

This is \cref{nab:thm:periodic} for $E=K_n=\Q(\zeta_n,2^{1/n})$, whose group is the
direct product of the affine groups (proof of \cref{grp:thm:artin2}): the $V_\ell$ are
distinct nonlinear irreducible characters, so $m_q=0$ and the minimum is
$\sum_\ell b_\ell^2$. If $\ell\mid q$ for all $\ell\in S$, \eqref{grp:eq:kummer-ind} gives
$\sum_\ell b_\ell\,\ell\,g_\ell=Z+\sum_\ell b_\ell I_\ell$ with $\sum_\ell b_\ell I_\ell$ a
function of $p\bmod q$.

\subsubsection{Consequences for Artin's conjecture}

\begin{theorem}[Prime indices suffice for the density; \stR]\label{grp:thm:artin2}
Suppose that the Dedekind zeta function of $K_\ell=\Q(\zeta_\ell,2^{1/\ell})$
satisfies the Riemann hypothesis for every prime $\ell$. Then
\[
 \#\set{p\le x:2\text{ is a primitive root modulo }p}\sim A\,\frac x{\log x},
 \qquad
 A=\prod_{\ell\ \mathrm{prime}}\Bigl(1-\frac1{\ell(\ell-1)}\Bigr)>0 .
\]
In particular, under Dirichlet GRH, the inequalities \eqref{grp:eq:nonres}
for every $\beta>1/2$ on a collection of finite blocks covering the odd primes
imply this asymptotic. No hypothesis for squarefree composite indices is
needed for this conclusion.
\end{theorem}

\begin{proof}
For $p\ne2$, $2$ fails to be a primitive root if and only if some prime
$\ell\mid p-1$ satisfies $2^{(p-1)/\ell}=1$ in $\F_p$, that is, $p$ splits
completely in $K_\ell$. For squarefree $n$ the field
$K_n=\Q(\zeta_n,2^{1/n})$ is the compositum of the $K_\ell$, $\ell\mid n$: one
inclusion is by taking powers and the other by B\'ezout's identity applied to
the integers $n/\ell$, with compatible positive real roots. The field
$\Q(\zeta_n)$ is unramified at $2$ (for even squarefree $n$ it equals
$\Q(\zeta_{n/2})$), so $X^n-2$ is Eisenstein at every prime over $2$ and
$[K_n:\Q]=n\varphi(n)$. For fixed $Y$ let $N_Y(x)$ count the primes $p\le x$
passing every test with $\ell\le Y$. Finite inclusion--exclusion and the
unconditional Chebotarev theorem for the finitely many $K_n$ give
\[
 \lim_{x\to\infty}\frac{N_Y(x)}{\pi(x)}
 =\prod_{\ell\le Y}\Bigl(1-\frac1{\ell(\ell-1)}\Bigr).
\]
We bound the tail uniformly, using the Riemann hypothesis only for the
prime-index fields. Put $y=\sqrt x/\log^4x$ and $z=\sqrt x\log^4x$. For a
Galois field $K$ whose zeta function satisfies the Riemann hypothesis, the
effective split-prime estimate \cite{Moree2012Artin,FanPollack2025Counting} is
\[
 \pi_{\mathrm{split}}(x;K)=\frac{\Li(x)}{[K:\Q]}
 +O\Bigl(\sqrt x\Bigl(\log x+\frac{\log\abs{D_K}}{[K:\Q]}\Bigr)\Bigr)
\]
with an absolute constant. The tower and the polynomial $X^\ell-2$ give
\[
 \abs{D_{K_\ell}}\le\abs{D_{F_\ell}}^\ell\bigl(\ell^\ell2^{\ell-1}\bigr)^{\ell-1},
 \qquad\abs{D_{F_\ell}}=\ell^{\ell-2},
\]
for odd $\ell$; hence $\log\abs{D_{K_\ell}}/[K_\ell:\Q]\ll\log(2\ell)$ and the
estimate is $\Li(x)/(\ell(\ell-1))+O(\sqrt x\log x)$ uniformly for
$\ell\le x$. For $Y<\ell\le y$, summing these bounds gives
\[
 \ll\Li(x)\sum_{\ell>Y}\frac1{\ell(\ell-1)}+y\sqrt x\log x
 \ll\frac x{Y\log x}+\frac x{\log^3x}.
\]
For $y<\ell\le z$ disregard the power condition and use the Brun--Titchmarsh
inequality for $p\equiv1\pmod\ell$: since $\log(x/\ell)\gg\log x$, the bound is
$O\bigl((x/\log x)\sum_{y<\ell\le z}1/(\ell-1)\bigr)$, and the sum over primes
is $O(\log\log x/\log x)$; this is $o(x/\log x)$. For $\ell>z$, a prime failing
the test divides $2^m-1$ with $m=(p-1)/\ell\le x/z$; the product of $2^m-1$
over $1\le m\le x/z$ has logarithm $O((x/z)^2)$, so it has $O(x/\log^8x)$
distinct prime divisors. Consequently $N_Y(x)$ differs from the desired count
by at most $O(x/(Y\log x))+o(x/\log x)$. Let $x\to\infty$ and then
$Y\to\infty$; the product converges to a positive value. The field $K_2$ is
quadratic and covered by Dirichlet GRH, and \cref{grp:thm:block} handles the
odd primes of the stated blocks.
\end{proof}

The theorem gives the density asymptotic, not the sharper error term of
Hooley's theorem, which uses the Riemann hypothesis for all squarefree
indices. For a general base the corresponding statement is the following.

\begin{corollary}[Reduction of Artin's conjecture; \stR]\label[corollary]{grp:cor:artin-general}
Assume Dirichlet GRH, and let $a$ be a nonzero integer, not $-1$ and not a
square. If the inequalities of \cref{grp:cor:kummer-general} hold for every
squarefree $m\ge2$, then Artin's primitive-root conjecture holds for $a$, with
Hooley's asymptotic density.
\end{corollary}

\begin{proof}
Each inequality implies the Riemann hypothesis for $\zeta_{E_m}$; the case
$m=1$ is the Riemann hypothesis for $\zeta$. Hooley's theorem
\cite{Hooley1967Artin,Moree2012Artin} uses exactly these hypotheses. The
detector is a fixed-field statement, but its conclusion is the exact zero
statement; Hooley's uniform estimates follow from the standard analytic theory
under that zero statement, not from uniformity of the constants
$C_{a,m,\eps}$.
\end{proof}

\subsubsection{The arithmetic obligation for a block}
Fix $S$, $h$, $b_\ell$ as in \cref{grp:thm:block} and let
$T(v)=-\sum_{\ell\in S}b_\ell\,\ell\,\overline R_{f_\ell}(v)$. Let $U$ be a
reference error satisfying the hypotheses of the window theorem
\cref{ana:thm:window}, for example a centred ramp error of a congruence class,
and let $C_0$ be the constant of that theorem for the pair $T,U$.

\begin{proposition}[Exact arithmetic obligation; \stR]\label[proposition]{grp:prop:bridge}
Assume Dirichlet GRH. Suppose there are fixed $A\ge1$ and $k>0$ such that for
each sufficiently small $\eps>0$, all large $v$, every $u\in[t,C_0t]$ with
$t=Av$, and every $0\le X\le e^t$ fixed independently of $u$, there are
admissible data giving, with a number $m\ge e^{v/2}$ of lifts and a positive
normalization $D_0$,
\begin{equation}\label{grp:eq:newstar}
 m\,D_0\bigl(T(v)-T(u)+X\,U(u)\bigr)\le C_{S,\eps}\,D_0\,e^{k\eps v}.
\end{equation}
Then $\zeta_{K_\ell}$ satisfies the Riemann hypothesis for every $\ell\in S$.
If this is available for blocks covering every odd prime, the asymptotic of
\cref{grp:thm:artin2} follows.
\end{proposition}

\begin{proof}
The transforms of $T$ and $U$ satisfy the hypotheses of \cref{ana:thm:window}
by \cref{grp:prop:interface}(a) and \cref{grp:lem:moving}: they are continuous,
of polynomial growth, and their transforms are holomorphic on a neighbourhood
of $\Re s\ge1$. Choose $u\in[t,C_0t]$ with $\abs{T(u)-XU(u)}\le e^{-t}$ and
divide \eqref{grp:eq:newstar} by $mD_0$ to get
$T(v)\le C_{S,\eps}e^{(-1/2+k\eps)v}+e^{-Av}$. For $\beta>1/2$ choose $\eps$
with $k\eps<\beta-1/2$ and apply \cref{grp:thm:block}.
\end{proof}

The inequality \eqref{grp:eq:newstar} models the source inequality
$m\cdot\mathrm{remainder}\le(\log b)^2\cdot\mathrm{budget}$ and its
cancellation at two scales; the data may depend on $u$, but $X$ must be fixed
at the first scale. It is not established. A proof would have to supply: an
integral lattice or jet-divisibility statistic whose local contributions
discriminate $2$ as an $\ell$th power for $p\equiv1\pmod\ell$, with the
orientation of \eqref{grp:eq:nonres}, which by \cref{grp:cor:periodic} no
periodic weight achieves; sufficiently many independent short forms and a
lower bound for the covolume of the same lattice, with matching main heights
and with $m\ge e^{v/2}$ forms; a justification of the averaging, since
unrelated multiplicative covolume inequalities may not be averaged before
taking logarithms and normalizing; and uniform selection of twists if the
congruence coefficients vary with the second scale, since a vector of moments
supplies only a common convex combination and not simultaneous pointwise
estimates.

\subsubsection{The case \texorpdfstring{$\ell=3$}{l=3}: all class weights}
Let $E=\Q(\sqrt[3]2)=B_3$ and $M=E(\zeta_3)=K_3$. The polynomial $X^3-2$ is Eisenstein at
$2$ and has discriminant $-108$; its splitting field $M$ has group $S_3$ and quadratic
subfield $\Q(\sqrt{-3})$. For $p>3$ let $g_p$ be the Frobenius class and $r_p$ the number
of roots of $X^3-2$ modulo $p$, and let $\tau=V_3$ be the character of the standard
representation, so that $\tau(g_p)=r_p-1$.
\begin{center}
\begin{tabular}{lrrrr}
\toprule
Factorization mod $p$ & Class size & $r_p$ & $\hecsgn(g_p)$ & $\tau(g_p)$\\
\midrule
Three linear factors & 1 & 3 & $1$ & $2$\\
One linear and one quadratic factor & 3 & 1 & $-1$ & $0$\\
One irreducible cubic & 2 & 0 & $1$ & $-1$\\
\bottomrule
\end{tabular}
\end{center}
The permutation character on the roots is $1+\tau$, so the standard factor is
$F=\zeta_E/\zeta=L_3$ of \eqref{grp:eq:kummer-fact}; it is entire, and $F(1)$ is finite
and positive.

\begin{example}[Three splitting types; \stR]\label[example]{hec:thm:cuberoot}
Assume \cref{src:thm:dgrh}. Let $w=(u,v,z)$ be real weights on the three rows of the
table, in order, and put
\[
 A=\frac{u+3v+2z}6,\qquad B=\frac{u-3v+2z}6,\qquad C=\frac{u-z}3,
\]
so that $w=A+B\,\hecsgn+C\tau$. Let
\[
 P_w(x)=\sum_{3<p\le x}w(g_p)\log p\,\Bigl(\frac1p-\frac1x\Bigr),\qquad
 R_w(x)=P_w(x)-A\log x-c_w,
\]
where the constant $c_w$ is fixed by $R_w(x)\to0$.
\begin{enumerate}[label=\textup{(\arabic*)}]
\item If $C=0$, then $R_w(x)=O_w(x^{-1/2}\log^2(2x))$.
\item If $C\ne0$, then RH for $\zeta_E$ is equivalent to the following one-sided estimate:
for every $\eps>0$ there is $K_{w,\eps}>0$ such that, for all sufficiently large $x$,
\[
 \hecsgn(C)\,R_w(x)\le K_{w,\eps}x^{-1/2+\eps}.
\]
Under RH for $\zeta_E$ the two-sided bound of \textup{(1)} holds for these weights as well.
\end{enumerate}
\end{example}

Here $S=\set{2,3}$ and $R_w=AR_\triv+BR_{\hecsgn}+CR_\tau$ with the remainders of
\cref{grp:prop:interface}(a), and $R_w(x)=-\int_x^\infty(\theta_w(t)-At)t^{-2}\,\dd t$.
If $C=0$, then $w$ is a combination of $\triv$ and $\triv+\hecsgn$; the quadratic field
gives $1+\hecsgn(g_p)$ degree-one primes above $p>3$, and the terms of higher residue
degree and higher powers add $O(\sqrt t\log^2(2t))$. Under RH for the respective Dedekind
zeta function the prime ideal theorem has error $O(\sqrt t\log^2(2t))$
\cite{GrenieMolteni2016Explicit}; for the quadratic field this hypothesis is supplied by
$\zeta_{\Q(\sqrt{-3})}=\zeta L(s,\chi_{-3})$ and \cref{src:thm:dgrh}. Hence
$\theta_w(t)-At=O_w(\sqrt t\log^2(2t))$, and the integral proves (1). By linearity the estimate in
(2) is equivalent to $R_\tau(x)\le K'_\eps x^{-1/2+\eps}$. The character
$\tau=\Ind_{A_3}^{S_3}\psi$, with $\psi$ nontrivial on $A_3$, is real-valued, and
$L(s,\tau)=F$; the case $a=0$ of \cref{grp:cor:detector} shows that this bound for every
$\eps$ is equivalent to $F\ne0$ in $\Re s>1/2$, that is, to RH for $\zeta_E$. Then
$\zeta_M=\zeta\,L(s,\chi_{-3})F^2$, by \eqref{grp:eq:kummer-fact} for $\ell=3$, also
satisfies the Riemann hypothesis, and the last assertion of \cref{grp:prop:interface}(d)
gives the two-sided bound for $R_\tau$.

Thus an upper bound for the primes at which $X^3-2$ splits completely suffices, and so
does a lower bound for the primes at which it is irreducible. For $w=(2,0,-1)$ the
logarithmic main term vanishes; this choice measures exactly the part missed by the two
Dirichlet characters in the table, and by \cref{nab:cor:nonlinear} its least mean-square
distance from a function of $p$ modulo a fixed modulus is $1$. The example identifies the
first missing instance of Hooley's family, with its sign and arithmetic weight; the full
conjecture would require the whole family, with the uniformity used in the sieve
\cite{FanPollack2025Counting}.

\subsection{A nonabelian arithmetic realization: the real \texorpdfstring{$S_3$}{S3} field of discriminant \texorpdfstring{$229$}{229}}
\label{grp:sec:s3}
{\rightskip=0pt plus 1fil\relax
\sources{grh_transfer_synthesis_2026-11-25; grh_nonabelian_transfer_2026-11-25}
}

The detectors above are analytic targets. This subsection concerns the other
side of the transfer: a specific nonabelian field which can be adjoined to the
CM field of the source construction (\cref{src:sec:mechanism}) at no cost in
ramification, and for which the local theta mechanism provably carries
nonabelian Frobenius information. These are arithmetic inputs. They do not
supply a global norm inequality, and no nonabelian Riemann hypothesis follows
from them.

\subsubsection{The field and its compatible embeddings}
Let
\[
 f(X)=X^3-4X+1,\qquad F_0=\Q(\alpha),\ f(\alpha)=0,\qquad
 F=\text{the splitting field of }f .
\]
By \cref{spn:ex:cubic229} and the discussion after it, $f$ is irreducible with three
real roots and prime discriminant $229$, $\Gal(F/\Q)\simeq S_3$, and only $229$ ramifies
in $F$, tamely, with inertia generated by a transposition. Hence $F$ is totally real and
$F=F_0(\sqrt{229})$. Put
$k=\Q(\sqrt{229})$, the quadratic subfield; it has conductor $229$.

Let $K=\Q(\sqrt{-d})$ be the imaginary quadratic field of the source
construction and let $q$ be a cyclotomic level with $229\mid q$. Put
\[
 E_0=K\Q(\zeta_q),\qquad E=FE_0 .
\]
Since $F\cap\Q(\zeta_q)$ is an abelian subextension of the $S_3$-extension
$F$, it lies in $k$, and $k\subseteq\Q(\zeta_{229})$; so
\begin{equation}\label{grp:eq:s3-fibre}
 F\cap\Q(\zeta_q)=k,\qquad
 \Gal(F\Q(\zeta_q)/\Q)\simeq
 \set{(g,u)\in S_3\times(\Z/q\Z)^\times:\grpsgn(g)=\chi_{229}(u)},
\end{equation}
a group of order $3\varphi(q)$. The source's splitting condition makes every
prime dividing $q$ split in $K$, and $K$ ramifies at $d$, where
$F\Q(\zeta_q)$ is unramified; so $K$ is disjoint from $F\Q(\zeta_q)$, the
induced CM type of $E$ has $3\varphi(q)$ embeddings, and $\Gal(E/K)$ is the
fibre product \eqref{grp:eq:s3-fibre}. Replacing it by a direct product would
give the wrong dimension and the wrong prime-splitting density.

\begin{theorem}[Unramified cubic base change; \stU]\label{grp:thm:unramified}
The extension $F/k$ is cyclic of degree three and unramified at every finite
prime. With $E_0$ and $E$ as above, $[E:E_0]=3$, the extension $E/E_0$ is
unramified at every finite prime, and
\begin{equation}\label{grp:eq:different}
 \grpDiff_{E/K}=\grpDiff_{E_0/K}\OO_E,\qquad
 \grpDiff_{E/\Q}=\grpDiff_{E_0/\Q}\OO_E,\qquad
 \grprd(E)=\grprd(E_0),
\end{equation}
where $\grprd$ denotes the root discriminant.
\end{theorem}

\begin{proof}
The subgroup fixing $k$ is $A_3$, so $F/k$ is cyclic cubic. Away from $229$
the splitting field is unramified, and at $229$ the inertia group is generated by a
transposition (\cref{spn:ex:cubic229} and the discussion after it); its intersection with $A_3$, which is the
inertia group for $F/k$, is trivial. So the relative different of $F/k$ is the
unit ideal. The field $E_0/\Q$ is abelian and contains $k$, hence
$F\cap E_0=k$ and $[E:E_0]=3$. An unramified extension of local fields remains
an unramified finite \'etale algebra after base change; applying this at all
finite places of $k$ gives $\grpDiff_{E/E_0}=\OO_E$. Transitivity of the
different gives the first two identities, and the relative discriminant
formula gives $\abs{\disc(E)}=\abs{\disc(E_0)}^3$ with $[E:\Q]=3[E_0:\Q]$.
\end{proof}

Thus a principal generator of the old different remains a generator after the
extension. This is special to the unramified benchmark; it does not say that
an arbitrary fractional ideal becomes principal or that adjoining an arbitrary
nonabelian field has the same property.

\subsubsection{The lattice, the Weil data and the unit witnesses}
From now on $q$ is an admissible tower modulus of the source with $229$ in its
fixed base modulus, and $A_0$ is the source's fractional Weil lattice in
$E_0$, whose local construction gives
$A_0\overline{A_0}=\grpDiff_{E_0/K}^{-1}$: inverse-different denominators are
permitted on one side of each selected conjugate pair of primes
(\code{finiteImage_mem_iff} in \code{New/Derivatives/Vacuum/ShortVectors}).
For the positive definite Hermitian pairing
$h(x,y)=\Tr_{E/K}(\overline xy)$ the dual of a fractional ideal $A$ is
\begin{equation}\label{grp:eq:hdual}
 A^\#=\grpDiff_{E/K}^{-1}\overline A^{-1},
\end{equation}
since $h(z,A)\subseteq\OO_K$ is equivalent to
$\overline zA\subseteq\grpDiff_{E/K}^{-1}$ and the different is stable under
conjugation \cite{Milne2020Algebraic}.

\begin{corollary}[The corrected lattice; \stU]\label[corollary]{grp:cor:lattice}
The ideal $A=A_0\OO_E$ satisfies $A\overline A=\grpDiff_{E/K}^{-1}$ and
$A^\#=A$. It is a full projective $\OO_K$-lattice in $E$ with a positive
definite Hermitian trace pairing.
\end{corollary}

\begin{proof}
Extend $A_0\overline{A_0}=\grpDiff_{E_0/K}^{-1}$ to $\OO_E$ and apply
\cref{grp:thm:unramified}; then \eqref{grp:eq:hdual} gives self-duality. A
fractional ideal is finitely generated and torsion-free over the Dedekind
domain $\OO_K$, hence projective. Under the induced CM type, the trace of
$\overline xx$ is a sum of squared absolute values.
\end{proof}

More generally, if every rational prime below $\grpDiff_{E/K}$ splits in $K$,
one may choose one of its two primes in $K$, put the valuation
$-v_{\mathfrak P}(\grpDiff_{E/K})$ in $A$ above the chosen prime and zero above
its conjugate; then $A\overline A=\grpDiff_{E/K}^{-1}$ prime by prime, with no
principality of the different required.

\begin{remark}[Weil data and unit witnesses]\label[remark]{grp:rem:weil}
The source library contains the statement
\code{New.nonempty_weilData_of_not_dvd_differentIdeal} (in
\code{Lib/ClassField/Grunwald}): a CM extension with compatible CM types
inherits Weil data when no finite prime fixed by conjugation divides the
relative different. \Cref{grp:thm:unramified} verifies its hypothesis in the
strongest form, and here the old finite idele can be kept by extension. The
other structural conditions transport: the root of unity of order $5040$ lies
in the base, the purely imaginary parameter has the same positive imaginary
values on the induced CM type, its trace-one algebraic integer remains
integral, and the splitting character is the original one composed with the
idele norm from $E$ to $E_0$. This is an application of an existing theorem of
the library, not a compiled instantiation.

The source frame construction also needs, for each $p$-adic embedding, a frame
member all of whose normalized CM-coordinate coefficients are units; in the
source the two possible witnesses are $1$ and
$\zeta_q\cdot\code{dualVector}(\varphi(q)-1)$, and the property is the norm
assertion $\abs{\iota_p(c\,\overline{\tau(y)})}_p=1$ for every CM embedding
$\tau$, with the period normalization $c$. Every new CM embedding restricts to
an old one, and the finite-idele valuations and the root discriminant are
unchanged, so the same norm assertions hold for the two old witnesses in $A$.
The vacuum root, defined in the source by a principal complex root, may change
by a root of unity under the degree-three base change, the quantity being
cubed and the root degree tripled; the statement transports as an equality of
absolute values at all places, or literally if the inherited old root is
chosen as the new normalization.
\end{remark}

The linear-algebra part of the frame input is treated in \cref{sec:spn}: $F$
has an integral tight frame of bound $229$ (\cref{spn:thm:optimal229}), which is optimal
because every ramified prime divides the bound of an integral frame, the elements
$v_i\zeta_q^j$ give a frame of bound $229q$ (\cref{spn:cor:tensor}) on the compatible embeddings of
$F\Q(\zeta_q)$ and on the induced CM type of $E$, and the self-dual projective
lattice $A$ admits frames containing prescribed vectors at every sufficiently
large integer bound. The frame of bound $229q$, divided by its bound and with
multiplicities multiplied by the bound, is a Parseval resolution which is
integral above every prime inert in $K$, since every prime factor of $229q$
splits in $K$.

\subsubsection{The local nonabelian partner and the splitting selector}

\begin{proposition}[Frobenius partner at inert primes; \stU]\label[proposition]{grp:prop:partner}
Let $p$ be inert in $K$, let $\iota:E\hookrightarrow\Cp$ select a prime
$\mathfrak p$ of $E$, and let $h\in\Gal(E/\Q)$ be the arithmetic Frobenius at
$\mathfrak p$. Let $c$ be complex conjugation, and identify the induced CM
type with the embeddings restricting to the chosen embedding of $K$. For odd
$e\ge1$ put $\operatorname{pr}_{p,e}(\tau)=c\,h^e\tau$. This is again a member
of the CM type, and for every $a\in\OO_E$
\begin{equation}\label{grp:eq:partner}
 \iota(\tau(a))^{p^e}\equiv
 \iota\bigl(\overline{\operatorname{pr}_{p,e}(\tau)(a)}\bigr)
 \pmod{p\,\OO_{E_{\mathfrak p}}}.
\end{equation}
On the components $F$ and $\Q(\zeta_q)$ the partner acts by $\Frob_{F,p}^e$
and by $-p^e$ respectively.
\end{proposition}

\begin{proof}
The prime $p$ is unramified in every factor of $E$, so
$\iota(h(b))\equiv\iota(b)^p\pmod{p\,\OO_{E_{\mathfrak p}}}$ for
$b\in\OO_E$. Iterating with $b=\tau(a)$ gives
$\iota(h^e\tau(a))\equiv\iota(\tau(a))^{p^e}$, and the conjugate of
$ch^e\tau(a)$ is $h^e\tau(a)$. Both $c$ and $h^e$ conjugate $K$, since $p$ is
inert and $e$ is odd, so their product fixes its chosen embedding. Complex
conjugation fixes the totally real field $F$ and inverts $\zeta_q$.
\end{proof}

Unramifiedness is essential: it makes the Frobenius congruence a divisibility
by $p$, the strength used in the source; a congruence modulo the maximal ideal
of a ramified completion would be weaker. Changing the $p$-adic embedding
conjugates Frobenius, and the selector below is invariant under this
ambiguity. The source's separate computation at the prime $p=d$ ramified in
$K$ is not covered.

To make the local connection explicit, let $u_j\in E^n$ be a Parseval
resolution, integral at every place above $p$, with multiplicities $a_j$, and
put $L_j(z)=\sum_{\tau,r}\overline{\tau(u_{j,r})}\,z_{\tau,r}$. All
coefficients lie in $\OO_E\otimes_\Z\Z_{(p)}$. Applying \eqref{grp:eq:partner}
to the conjugated coefficients and using the Parseval identity gives
\begin{equation}\label{grp:eq:coupling}
 \sum_ja_jL_j(z)^{p^e}L_j(z)\equiv
 \sum_{\tau,r}z_{\tau,r}^{p^e}\,z_{\operatorname{pr}_{p,e}(\tau),r}\pmod p ;
\end{equation}
for instance the coefficient of $z_{\tau,r}^{p^e}z_{\sigma,s}$ on the left is
congruent to
$\sum_ja_j\operatorname{pr}_{p,e}(\tau)(u_{j,r})\overline{\sigma(u_{j,s})}$,
the required Kronecker delta. The generic source lemma
\code{dvd_sum_subst_sub_couplingSum} (in
\code{New/Derivatives/FormalGroup/ThetaProducts}) accepts such a coupling
congruence for an arbitrary partner map, and \code{HasNormalForm} imposes no
cyclotomic restriction on that map; \cref{grp:prop:partner} and
\eqref{grp:eq:coupling} replace the arithmetic input of the specialized
cyclotomic bridge at inert primes. The coefficients of the linear forms carry
nonabelian information although the one-variable theta function still comes
from a CM elliptic curve over $K$.

\begin{lemma}[Coset correlation; \stU]\label[lemma]{grp:lem:coset}
Let $G$ be a finite group, $H\le G$, $t\in G$, $S\subseteq H$ and $B=tS$. Then
for every $g\in G$
\[
 \#\set{a\in H:ga\in B}=\abs B\,\grpone_{gH=tH}.
\]
If instead $\pi:G\twoheadrightarrow Q$ is a quotient and
$A=\pi^{-1}(A_0)$, $B=\pi^{-1}(B_0)$, then
$\#\set{a\in A:ga\in B}=\abs{\ker\pi}\,\#\set{a_0\in A_0:\pi(g)a_0\in B_0}$.
\end{lemma}

\begin{proof}
Left multiplication by $g$ is a bijection from $H$ to $gH$. If $gH=tH$ its
inverse gives one preimage for each element of $B$; otherwise the two cosets
are disjoint. The second statement is a count in the fibres of $\pi$.
\end{proof}

Let $q_0\mid q$ be the base modulus, $B_0=F\Q(\zeta_{q_0})$, and
\[
 G=\Gal(E/K),\qquad H=\Gal(E/KB_0)=\set{(1,u):u\equiv1\pmod{q_0}},
 \qquad t=(1,-1),
\]
so that $[G:H]=3\varphi(q_0)$ and $H$ is central; its cyclotomic tower
coordinates permit a sparse $S\subseteq H$. For the partner $g=ch^e$ of an
inert prime, the condition $gH=tH$ of \cref{grp:lem:coset} is exactly
\begin{equation}\label{grp:eq:selector}
 \Frob_{B_0,p}^{\,e}=1 .
\end{equation}
For $e=1$ this is complete splitting in $B_0$, of density
$1/(3\varphi(q_0))$, and the formula retains the prime-power information for
higher $e$. This is the group-theoretic replacement for the place count
\code{card_filter_mem_bPlaces} of the source. By the second statement of the
lemma, selectors pulled back from the cyclotomic quotient cannot detect
Frobenius in the kernel; the selector $H,tS$ avoids this.

\begin{example}\label[example]{grp:ex:s3-primes}
Take $q_0=4231920=18480\cdot229$ and $d=167160839$. Then $d$ is prime,
$d\equiv7\pmod8$, $q_0\mid2(d+1)$ and $q_0\nmid d+1$. The primes
$p_1=38087281$ and $p_2=59246881$ are both $1$ modulo $q_0$ and inert in $K$,
yet
\[
 f(X)\equiv(X+16188066)(X-13084869)(X-3103197)\pmod{p_1},
\]
whereas $f$ is irreducible modulo $p_2$, since $\gcd(f,X^{p_2}-X)=1$. The
Frobenius classes in $F$ are the identity and a $3$-cycle. A congruence
$p\equiv1\pmod N$ does not force splitting in $F$, even when $N$ is divisible
by its discriminant. The verification is recorded in\\
\artifact{grh_transfer_synthesis_2026-11-25/validation/nonabelian_frame.json}.
\end{example}

The prime coefficient suggested by this local realization is
\[
 \log p\,(1-\chi_K(p)^e)\,\grpone_{\Frob_{B_0,p}^{e}=1},
\]
which concerns nonabelian prime splitting and so contains information absent
from the Dirichlet congruence mask of the source. The source's main-term
cancellation compares remainders at two cutoffs, so a class-specific constant
cancels; keeping the elliptic CM period from $K$ is therefore not by itself an
obstruction to this signal.

\subsubsection{The first analytic target and the quantitative budget}
Let $\varrho$ be the two-dimensional irreducible representation of $S_3$. Then
\begin{equation}\label{grp:eq:zetaF}
 \zeta_F(s)=\zeta(s)L(s,\chi_{229})L(s,\varrho)^2,
\end{equation}
and $L(s,\varrho)$ is the Hecke $L$-function over $k$ of a cubic character
attached to the unramified extension $F/k$; it is entire by Hecke's theorem,
and this is not a new case of Artin holomorphy. Write
$-\zeta_F'/\zeta_F(s)=\sum_na_F(n)n^{-s}$ and
$\zeta_F(s)=\kappa_F/(s-1)+\kappa_F\gamma_F+O(s-1)$. Then
$0\le a_F(n)\le6\Lambda(n)$, and at an unramified prime
$a_F(p^j)=6\log p\,\grpone_{\Frob_p^j=1}$. Define
\begin{equation}\label{grp:eq:targetR}
 R_F(x)=\sum_{n\le x}a_F(n)\Bigl(\frac1n-\frac1x\Bigr)-\log x+1+\gamma_F .
\end{equation}

\begin{proposition}[A sufficient final estimate; \stR]\label[proposition]{grp:prop:s3-target}
If for every $\eps>0$ and every sufficiently large dyadic $Y$
\begin{equation}\label{grp:eq:target-mean}
 \int_Y^{2Y}\grppos{R_F(x)}\,\dd x\ll_\eps Y^{1/2+\eps},
\end{equation}
then every nontrivial zero of $\zeta_F$ lies on $\Re s=1/2$, and the same
holds for the Hecke factor $L(s,\varrho)$.
\end{proposition}

\begin{proof}
The coefficient bound gives $R_F(x)=O(\log^2x)$, and termwise integration
gives, for $\Re s>1$,
\[
 \int_1^\infty R_F(x)x^{-s}\,\dd x
 =-\frac{\zeta_F'(s)/\zeta_F(s)}{s(s-1)}-\frac1{(s-1)^2}+\frac{1+\gamma_F}{s-1}.
\]
The singularity at $1$ is removable, and a zero $\omega$ of multiplicity
$m_\omega$ in $\Re s>0$ gives the residue $m_\omega/(\omega(1-\omega))$,
positive at real zeros in $(0,1)$. Summing \eqref{grp:eq:target-mean} over
dyadic intervals shows that the Mellin integral of the positive part converges
for $\Re s>1/2$. The positive-part form of the one-sided test
(\cref{ana:thm:onesided}, see \cref{grp:rem:variants}) then shows that the
transform has no pole in $\Re s>1/2$, so $\zeta_F$ has no zero there; the
functional equation gives the critical line. The other factors in
\eqref{grp:eq:zetaF} have no poles in the open critical strip, so an off-line
zero of $L(s,\varrho)$ would be a zero of $\zeta_F$.
\end{proof}

The local theta construction naturally produces first a difference with an
auxiliary quadratic twist, as in the source; reaching
\eqref{grp:eq:target-mean} directly, or proving the analogue of the source's
twist selection with cancellation excluded, is a separate task.

The source theorem \code{New.exists_length_shaOdd_le} (in
\code{Lib/Cyclotomic/OddSha}) requires the entire CM field to lie in a
cyclotomic field $\Q(\mu_N)$ and produces a nonzero integer
$X\le2^{4(N+t+3)^{27}}h_E$ whose valuations bound the lengths of the relevant
finite-level odd cocycle modules. A field containing $F$ lies in no cyclotomic
field, so this input has no direct analogue for $E$. Its required strength,
however, is much weaker than what the source proves: by
\cref{sel:prop:growth}, under the source's scale inequalities the odd-cocycle
contribution $t\log X$ to the final denominator budget
$\mathcal B=fa(\log b)^3/(mk)$ is negligible as soon as $\log X\le b$, and for
$\log b\asymp x^\eps$ even a uniform bound $\log X\le\exp(x^{o(1)})$ suffices
for this term, whereas the source bounds $\log X$ by $x^{2000}$. For $E$ this
relaxes the target; it does not bound the new modules, settle the coefficient
prime $3$, or account for the other denominator terms.

\begin{remark}[What the realization does not provide]\label[remark]{grp:rem:s3-missing}
The source's scale hierarchy takes $m\asymp\sqrt x$, $r\asymp x^5$, $t=5r$,
$\log b\asymp x^\eps$ and $f\ge\varphi(q_0)x^{11}$ for every small $\eps>0$,
with a tower modulus of polynomial growth in $x$. The following inputs of the
source argument are not available for $E$.
\begin{enumerate}[label=\textup{(\arabic*)}]
\item \emph{Uniform frames.} The frames of \cref{sec:spn} exist at every
sufficiently large bound, with a threshold depending on the field, the lattice
and the prescribed witnesses. A threshold growing like $\exp(f^c)$ would
exceed the permitted scale $b=\exp(x^\eps)$; bounds uniform in the growing
tower, for the reserve denominator, the signed coefficients, the vector sizes
and the threshold, are required, as are frames of the prescribed principal
levels.
\item \emph{Odd cocycles.} A bound for the source's raw odd cocycle modules,
with their local conditions and coefficient-ring quantifiers, for the
noncyclotomic field $E$, of the strength permitted by
\cref{sel:prop:growth}. Since $E/E_0$ is cyclic cubic and unramified, away
from the coefficient prime $3$ restriction and corestriction identify the
invariant part of ordinary cohomology with the contribution of $E_0$; the two
other eigenspaces involve the nontrivial cubic characters of $E/E_0$ and are
not controlled by the invariant part. The comparison of this description with
the raw cocycle modules is not proved.
\item \emph{Level primes.} Splitting and principal-ideal conditions for the
level primes inside the source's quantitative windows; by
\cref{grp:ex:s3-primes} congruence conditions do not suffice.
\item \emph{Local defects and spanning.} The crystalline-defect estimates of
the source use the cyclotomic nature of the field, and its spanning theorem
\code{isSpread_goodPoincare} is not inherited by tripling the number of
embeddings.
\item \emph{The final estimate.} A bound such as \eqref{grp:eq:target-mean},
or a matched version with an auxiliary twist and cancellation excluded.
\end{enumerate}
\end{remark}

\subsection{Open questions}\label{grp:sec:open}
{\rightskip=0pt plus 1fil\relax
\sources{grh_method_transfers_2026-11-25_1; grh_sharp_transfers_2026-11-25; grh_a6_transfer_2026-11-25; grh_sylow_transfer_2026-11-25; psl2_grh_transfer; grh_transfers_density_gaps_and_zero_tests; grh_spanning_and_s5_2026-11-26; grh-exponent-six-and-transfers; grh-transfer; grh_kummer_window_transfer_2026-11-25; grh_nonabelian_transfer_2026-11-25; grh_transfer_synthesis_2026-11-25; ProveIt_research_2026-11-27}
}

\begin{problem}[The prime inequalities]\label{grp:prob:inequalities}
Prove any one of the one-sided inequalities of this section. The most
economical instances are: the two-pattern dyadic inequality for the sextic
\eqref{grp:eq:a6-f} (\cref{grp:ex:a6-field}); the two-pattern test
\eqref{grp:eq:psl27-test} for Trinks' polynomial; the root-count inequality of
\cref{grp:cor:psl-root} for a field with group $\PSL_2(8)$ of degree nine; the
inequality \eqref{grp:eq:a5-bound} for $X^5+20X+16$; the quartic test of
\cref{grp:thm:s4}; and the first Kummer block $S=\set3$ of
\cref{grp:thm:block}, that is, a bound for the deficit of primes
$p\equiv1\pmod3$ modulo which $2$ is not a cube.
\end{problem}

\begin{question}[Arithmetic carriers of nonabelian weights]\label{grp:q:carrier}
Is there a height, lattice or jet-divisibility construction whose local
contribution at an unramified prime $p$ is a prescribed detector weight
$J(\Frob_p)\log p$, with matching main terms on both sides of a covolume
comparison and with errors preserving the square-root saving? Specific
instances: the permutation trace $P$ of $A_6$, which counts degree-one primes
of the field $L^{H_4}$ of degree $90$; a positive translate of $W_P/2$ for
$\PSL_2(7)$, over the field of degree $21$ and its two quadratic extensions;
the two-point and nonsplit-torus actions of $\PSL_2(q)$, where
$W=\pi_B+\pi_{T'}-\triv$ suggests pairing local valuations in two intermediate
fields; the weight $W_{21}$ of $A_5$ through the realizations of degrees
$5,6,12$; the cubic residuacity of $2$ among $p\equiv1\pmod3$ with the sign of
\eqref{grp:eq:nonres}; and the selector \eqref{grp:eq:selector} for the field
of discriminant $229$, for which the missing inputs are listed in
\cref{grp:rem:s3-missing}.
\end{question}

\begin{question}[Odd characteristic]\label{grp:q:odd}
For $q\equiv3\pmod4$, what is the least possible maximum degree of a spectrally
detecting family of subfields of a $\PSL_2(q)$-extension? For $q\equiv1\pmod4$
it is determined by \cref{grp:thm:odd-family}, and for every odd $q$ the
three fields of \cref{grp:prop:odd-three}, of largest degree $q(q-1)$, detect
zeros unconditionally. The least nonnegative dimension and the single-subfield
degree are settled by \cref{grp:thm:odd-gelfand,grp:thm:odd-rich}. Can Artin
holomorphy be removed from the equivalence for a single optimal field, such as
$L^{T_n}$ when $q$ is prime (\cref{grp:rem:odd-zeros})?
\end{question}

\begin{question}[Other groups]\label{grp:q:groups}
Which other finite simple groups admit a small nonnegative detector whose
$L$-function is holomorphic away from $1$ by induction from linear characters?
More generally, for a given group, minimize degree, support and the
multiplicity constant $k$ over identities
$kJ+b\,\triv=\Reg+\sum_ia_i\Ind_{H_i}^G\lambda_i$ with $a_i\ge0$. The pair
$W,Y$ for $A_6$ shows that these objectives conflict.
\end{question}

\begin{question}[Critical-line zeros]\label{grp:q:line}
Given one or more positive monomial detectors, can their local zero orders
constrain the Heilbronn character enough to exclude negative irreducible
coefficients on the critical line? The criteria of this section control the
open right half of the strip only.
\end{question}

\begin{question}[Kummer blocks]\label{grp:q:kummer}
Can a uniform family of analytic disks be obtained for increasing blocks in
\cref{grp:thm:block,grp:prop:bridge}, and can an arithmetic construction treat
a block aggregate more cheaply than the separate fields? A finite-block
statement with constants depending on the block suffices for the density
limit of \cref{grp:thm:artin2}; uniformity is needed for a quantitative error
term or for a single criterion with an infinite block. What replaces
\cref{grp:thm:artin2} for other bases and for Artin-type problems of rank
greater than one, where the abelian parts, the ramification and the
entanglement of the Kummer fields differ from the base $2$?
\end{question}

\begin{question}[The tame line]\label{grp:q:tame}
For $D=-12451$ exactly one of the six lines of \cref{grp:sec:six} gives a
quintic field tame at $5$; it corresponds to a cyclic quintic extension of the
quartic mirror field $H$ unramified at every finite prime. Can a theta-lift or
height construction estimate its splitting primes with the sign and
square-root accuracy of \cref{grp:thm:six}, first for this line alone?
\end{question}

\endgroup
